% !TeX program = xelatex
% Compact: xelatex COMP2610-cheatsheet.tex
% Large print: xelatex -jobname=COMP2610-worked-solutions '\def\LargePrint{1}% !TeX program = xelatex
% Compact: xelatex COMP2610-cheatsheet.tex
% Large print: xelatex -jobname=COMP2610-worked-solutions '\def\LargePrint{1}% !TeX program = xelatex
% Compact: xelatex COMP2610-cheatsheet.tex
% Large print: xelatex -jobname=COMP2610-worked-solutions '\def\LargePrint{1}% !TeX program = xelatex
% Compact: xelatex COMP2610-cheatsheet.tex
% Large print: xelatex -jobname=COMP2610-worked-solutions '\def\LargePrint{1}\input{COMP2610-cheatsheet.tex}'
\documentclass[a4paper]{article}
\usepackage[UTF8,scheme=plain,fontset=fandol]{ctex}
\usepackage{amsmath,amssymb,multicol,xcolor}
\ifdefined\LargePrint
  \usepackage[margin=14mm]{geometry}
  \setlength{\columnsep}{6mm}
  \newcommand{\bodyformat}{\fontsize{9.6}{13}\selectfont}
  \newcommand{\cols}{2}
  \setlength{\parskip}{2.7pt}
\else
  \usepackage[top=5.6mm,bottom=7.6mm,left=5.8mm,right=5.8mm]{geometry}
  \setlength{\columnsep}{3mm}
  \newcommand{\bodyformat}{\fontsize{6.65}{7.8}\selectfont}
  \newcommand{\cols}{4}
  \setlength{\parskip}{1.05pt}
\fi
\usepackage{fancyhdr}
\definecolor{ink}{HTML}{163D51}
\pagestyle{fancy}
\fancyhf{}
\fancyfoot[L]{\scriptsize COMP2610 / COMP6261}
\fancyfoot[R]{\scriptsize \thepage}
\renewcommand{\headrulewidth}{0pt}
\setlength{\footskip}{4.5mm}
\setlength{\parindent}{0pt}
\setlength{\multicolsep}{2pt}
\setlength{\emergencystretch}{1.3em}
\setlength{\columnseprule}{.14pt}
\setlength{\premulticols}{0pt}
\setlength{\postmulticols}{0pt}
\raggedcolumns
\newcommand{\topic}[1]{\par\addvspace{2.6pt}{\color{ink}\bfseries #1}\par\nobreak}
\newcommand{\problem}[2]{\par\addvspace{1.4pt}\textcolor{ink}{\textbf{#1}}\enspace\textbf{题}\enspace #2\par\nobreak}
\newcommand{\solve}{\noindent\textbf{解}\enspace}
\newcommand{\parta}[1]{\par\noindent\textbf{(#1)}\enspace}
\newcommand{\form}[1]{\par\noindent\textbf{#1}\enspace}
\newcommand{\E}{\mathbb E}
\newcommand{\V}{\operatorname{Var}}
\newcommand{\B}{h_2}
\newcommand{\bb}[1]{\texttt{#1}}
\newcommand{\code}{\mathrm{EOS}}
\begin{document}
{\color{ink}\fontsize{13}{15}\selectfont\bfseries COMP2610 / COMP6261 \quad L22 公式与 Tut 逐题解}\par
{\fontsize{7}{8.4}\selectfont Tut 1–12 共 46 道去重题。题干以中文保留条件和小问；Tut 5 Q4 与 Tut 7 Q2 共用一份解答。信息量对数以 2 为底。\par}
\begin{multicols}{\cols}
\bodyformat\raggedright
\topic{公式区}
\form{概率与 Bayes} $p(x,y)=p(x)p(y|x)$，$p(y)=\sum_xp(x,y)$，$p(x|y)=p(x,y)/p(y)$。独立时 $p(x,y)=p(x)p(y)$；条件独立在固定条件后用同一乘积形式。
\form{Expected Value} $\E g(X)=\sum_x p(x)g(x)$。期望可直接相加。
\form{Variance} $\V X=\E[X^2]-(\E X)^2$。独立变量相加时方差相加。
\form{Binomial Distribution} $K=\sum_{i=1}^n B_i$，$B_i$ 独立且为 $\mathrm{Bern}(p)$。$P(K=k)=\binom nk p^k(1-p)^{n-k}$，$\E K=np$，$\V K=np(1-p)$。理由：$B_i^2=B_i$，所以单次期望 $p$、方差 $p-p^2$。
\form{Markov's Inequality} $W\ge0,\ a>0$ 时，$P(W\ge a)\le\E W/a$。阈值以上每个结果至少贡献 $a$。
\form{Chebyshev's Inequality} $\mu=\E W,\ t>0$ 时，$P(|W-\mu|\ge t)\le\V(W)/t^2$。单侧阈值 $a>\mu$ 用 $t=a-\mu$。这是对 $(W-\mu)^2$ 用 Markov。
\form{Jensen / MLE} 凸 $f$ 满足 $f(\E X)\le\E f(X)$。i.i.d. 的 MLE 最大化 $\ln L(\theta)=\sum_i\ln p_\theta(x_i)$。
\form{Entropy} 自信息 $-\log p(x)$；$H(X)=-\sum_xp(x)\log p(x)$；$H(X,Y)=-\sum_{x,y}p(x,y)\log p(x,y)$。$H(Y|X)=\sum_xp(x)H(Y|X=x)$。$0\log0=0$。$H(X,Y)=H(X)+H(Y|X)$；$H(X)\le\log|\mathcal X|$。记 $\B(p)=-p\log p-(1-p)\log(1-p)$。
\form{互信息与 KL} $I(X;Y)=H(X)-H(X|Y)=H(Y)-H(Y|X)=D(p_{XY}\Vert p_Xp_Y)$；$D(p\Vert q)=\sum p\log(p/q)\ge0$。条件互信息是给定条件时互信息的加权平均。对 $X\to Y\to Z$，$I(X;Z)\le\min\{I(X;Y),I(Y;Z)\}$。
\form{Extended Ensemble} $X^N$ 把整块视作一个结果。i.i.d. 时 $p(x^N)=\prod_i p(x_i)$，$H(X^N)=NH(X)$。含 $k$ 个 1 的单串概率为 $p^k(1-p)^{N-k}$，同组有 $\binom Nk$ 串。
\form{AEP／典型集} $T_{N,\beta}=\{x^N:|-\log p(x^N)/N-H(X)|<\beta\}$，i.i.d. 时其概率趋 1。最小 $\delta$-充分集按单串概率降序取至累计概率 $\ge\delta$；$H_\delta=\log|S_\delta|$，定长需 $\lceil H_\delta\rceil$ 位。固定 $0<\delta<1$ 时 $H_\delta/N\to H(X)$。
\form{信源编码} $L=\sum_i p_i\ell_i$；二元 Kraft 和 $\sum_i2^{-\ell_i}\le1$。Huffman 每轮并两个最小权重。Shannon 取累计 $F_i=\sum_{j<i}p_j$ 的前 $\ell_i=\lceil-\log p_i\rceil$ 位；SFE 取 $F_i+p_i/2$ 的前 $\ell_i+1$ 位。$H\le L_{\rm opt}\le L_{\rm Shannon}<H+1$。
\form{算术编码} 当前 $[a,a+w)$，符号概率段 $[F,F+p)$，更新为 $[a+wF,a+w(F+p))$。含 EOS 的最终区间内取一个完整二进制网格区间。
\form{信道编码} $Q_{y,x}=p(y|x)$，$p_Y=Qp_X$；$H(Y|X)=\sum_xp(x)H(Q_{\cdot,x})$；$C=\max_{p_X}I(X;Y)$。$R=K/N$；$P_B=P(\hat{\mathbf s}\ne\mathbf s)$；$R<C$ 可使块错率趋零，$R>C$ 不可。

\topic{Tut 1 · 基础概率}
\problem{1.1}{等分转盘标 $1,2,3,4,5$，求 (a) 转到 4，(b) 转到偶数，(c) 转到质数的概率。}
\solve 五格各占 $1/5$，互斥格可加。(a) $\{4\}$：$1/5$。(b) $\{2,4\}$：$2/5$。(c) $\{2,3,5\}$：$3/5$。
\problem{1.2}{ACT 车牌是三字母接三数字。求随机车牌以 ACT 开头且以 7 结尾的概率。}
\solve 按各合法组合等可能且允许重复，总数 $26^3\cdot10^3$。固定前三字母和末位后，中间两位各有 10 种，共 $10^2$。概率 $10^2/(26^3\cdot10^3)=1/175760$。
\problem{1.3}{四条路的测速启用率 $(.4,.3,.2,.3)$，Arya 选路概率 $(.2,.1,.5,.2)$，且总超速。求 (a) 一早被罚，(b) 五早都不被罚。}
\solve (a) 对每条路先乘选路概率和该路处罚率，再相加，得 $.2(.4)+.1(.3)+.5(.2)+.2(.3)=.27$。(b) 每早不被罚概率 $1-.27=.73$。按各早独立且规则不变，五早全不罚为 $.73^5\approx.20731$。
\problem{1.4}{袋内 5 蓝、3 红、2 黄，抽三球。求红球少于 2 颗的概率，分别按 (a) 放回，(b) 不放回。}
\solve (a) 红概率 $.3$，非红 $.7$，三次独立。零红 $.7^3$；一红的出现位置有 3 个，概率 $3(.3)(.7)^2$；和为 $.784$。(b) 把无序三球子集当等可能结果。零红选三非红，有 $\binom73=35$ 种；一红选一红两非红，有 $\binom31\binom72=63$ 种。总数 $\binom{10}3=120$，故 $98/120=49/60$。
\problem{1.5}{两台录像机独立成功，概率分别为 $.7,.6$。求 (a) 零份，(b) 一份，(c) 两份录像的概率。}
\solve 失败率分别是 $.3,.4$。(a) 两失败：$.3(.4)=.12$。(b) 仅第一台或仅第二台，两情形互斥：$.7(.4)+.3(.6)=.46$。(c) 两成功：$.7(.6)=.42$。三项和为 1。

\topic{Tut 2 · 条件概率与 Bayes}
\problem{2.1}{红、蓝、绿箱被选概率为 $(.2,.2,.6)$，各箱苹果、橙子、青柠数分别为 $(3,4,3)$、$(1,1,0)$、$(3,3,4)$。选箱后均匀抽一果。求 (a) 苹果概率，(b) 给定橙子，来自绿箱的概率。}
\solve (a) 各箱苹果比例 $3/10,1/2,3/10$；乘箱概率后总和为 $.06+.10+.18=.34$。(b) 橙的三项联合概率为 $.2(4/10)=.08$、$.2(1/2)=.10$、$.6(3/10)=.18$。用绿箱项除总和，得 $.18/.36=1/2$。
\problem{2.2}{共 1000 次，行 $Y=0,1$、列 $X=0,1$ 的计数为 $(100,250)$、$(150,500)$。求 (a) $p(1,1)$，(b) $P(X=1)$，(c) $\E X$，(d) $P(Y=1|X=1)$，(e) $P(Y=1|X=0)$。(f) 若 $X\ne Y$ 时 $P(Z=1|X,Y)=.9$，相同则为 $.1$，求 $P(X=1,Y=1|Z=1)$。}
\solve (a) 单格 $500/1000=.5$。(b) $X=1$ 列和 $(250+500)/1000=.75$。(c) $0(.25)+1(.75)=.75$。(d) 固定 $X=1$，分母 750，得 $500/750=2/3$。(e) 固定 $X=0$，分母 250，得 $150/250=.6$。(f) 不同格合计 $400/1000=.4$。全概率 $P(Z=1)=.9(.4)+.1(.6)=.42$。目标格且 $Z=1$ 的联合概率是 $.5(.1)=.05$，所以后验 $=.05/.42=5/42$。
\problem{2.3}{独立抛公平硬币两次，每次正面放白球，反面放红球。随后从这两球中放回抽三次，均为红。求箱内两球都红的后验概率。}
\solve 令红球数 $K$。由两次独立抛币，$K=0,1,2$ 的先验为 $(1,2,1)/4$。固定 $K=k$ 后，放回抽每次红率 $k/2$，三红似然为 $(k/2)^3$。乘先验得未归一化后验 $(0,1/16,1/4)$。故 $P(K=2|\text{三红})=(1/4)/(1/16+1/4)=4/5$。
\problem{2.4}{(a) 从条件概率的定义证明 $p(x|y,z)=p(y|x,z)p(x|z)/p(y|z)$。(b) 证明给定 $Y$ 时仍可对任一变量 $X_i$ 求和边缘化。}
\solve (a) 在分母为正时，$p(x|y,z)=p(x,y|z)/p(y|z)$。分子按乘法规则为 $p(x|z)p(y|x,z)$，代回即得。(b) 令 $U$ 代表除 $X_i$ 外的变量。$p(u|y)=p(u,y)/p(y)=[\sum_x p(u,x,y)]/p(y)=\sum_x p(u,x|y)$。

\topic{Tut 3 · MLE 与熵}
\problem{3.1}{取值 $1,2,3$ 的概率为 $(\theta/2,\theta/2,1-\theta)$，$0\le\theta\le1$。i.i.d. 样本中各频数是 $n_1,n_2,n_3$。(a) 写似然。(b) 样本 $3,3,1,2,3,2,2,1,3,1$ 求 MLE。}
\solve (a) 把单次概率相乘，同类因子合并，得 $L=(\theta/2)^{n_1+n_2}(1-\theta)^{n_3}$。(b) 频数 $(3,3,4)$；$\ln L=6\ln(\theta/2)+4\ln(1-\theta)$。令导数 $6/\theta-4/(1-\theta)=0$，得 $6(1-\theta)=4\theta$，即 $\theta=.6$。导数在此点左右分别正负，两端似然为 0，故为最大值。
\problem{3.2}{联合表的行是 $Y=1,2,3,4$，列是 $X=1,2,3,4$，矩阵按 $1/16$ 为单位，四行为 $(0,0,2,2)$、$(2,1,1,0)$、$(2,2,0,0)$、$(0,1,1,2)$。(a) 证不独立。(b) 求 $H(X)$、$H(Y)$、$H(X|Y)$、$H(Y|X)$、$H(X,Y)$。}
\solve (a) 各行各列都合计 $4/16=1/4$。左上为 0，而独立时应为 $1/4\cdot1/4=1/16$。(b) 两个边缘均匀，两个熵都是 $4(1/4)\log 4=2$。每行除以 $1/4$，非零条件概率为 $(1/2,1/2)$ 或 $(1/2,1/4,1/4)$，熵分别为 $1$ 或 $1/2+1/2+1/2=1.5$。四行熵 $(1,1.5,1,1.5)$，均值 $1.25=H(X|Y)$。四列熵 $(1,1.5,1.5,1)$，均值 $1.25=H(Y|X)$。联合表有六个 $1/8$ 和四个 $1/16$，所以 $H(X,Y)=6(1/8)3+4(1/16)4=3.25$。
\problem{3.3}{均匀抽 52 张牌之一。$V,S$ 为点数和花色，人头牌是 J、Q、K。求 (a) 看到红 K 的自信息，(b) 已知是人头牌时看到 K 的条件自信息，(c) $H(S)$ 与 $H(V,S)$。}
\solve (a) 红 K 有两张，$-\log(2/52)=\log26$。(b) 12 张人头牌中四张 K，$-\log(4/12)=\log3$。(c) 花色各有 $13/52=1/4$，$H(S)=2$；一对 $(V,S)$ 确定一张牌，52 对等可能，$H(V,S)=\log52$。
\problem{3.4}{有限离散 $X$。比较 $H(X)$ 与 $H(Y)$，分别取 (a) $Y=2^X$，(b) $Y=\cos X$。}
\solve 确定映射给 $H(Y|X)=0$，两种顺序展开联合熵得 $H(X)=H(Y)+H(X|Y)$。(a) $X=\log_2Y$，故 $H(X|Y)=0$，两熵相等。(b) $H(X|Y)\ge0$，故 $H(Y)\le H(X)$。若 $X$ 等概率为 $\pm1$，余弦相同，熵从 1 变 0；若在支持集上可逆则取等。

\topic{Tut 4 · 熵、Jensen、互信息}
\problem{4.1}{几何变量 $Y$ 的概率为 $p(1-p)^{y-1}$，$y=1,2,\dots$。求期望和方差。}
\solve 设 $q=1-p$、$m=\E Y$、$s=\E Y^2$。首轮成功时 $Y=1$；失败时 $Y=1+Y'$，$Y'$ 同分布。故 $m=p+q(1+m)=1+qm$，得 $m=1/p$。平方后 $s=p+q(1+2m+s)=1+2qm+qs$，故 $s=(1+2q/p)/p=(2-p)/p^2$；方差 $s-m^2=(1-p)/p^2$。
\problem{4.2}{A、B 两队独立且各半胜率，先赢四局者获胜。$X$ 是完整胜负序列，$Y$ 是总局数。求 (a) $H(X)$，(b) $H(Y)$，(c) $H(Y|X)$，(d) $H(X|Y)$。}
\solve 选总冠军有两种；总长 $y$ 时第 $y$ 局必由冠军赢，前 $y-1$ 局选出其前三胜，故 $M_y=2\binom{y-1}3$。$y=4,5,6,7$ 的序列数为 $2,8,20,40$。每序列概率 $2^{-y}$，所以 $p_Y=(1/8,1/4,5/16,5/16)$。
\parta{a} 序列自信息为 $y$，故 $H(X)=4/8+5/4+6(5/16)+7(5/16)=5.8125$。
\parta{b} $H(Y)=(1/8)3+(1/4)2-2(5/16)\log(5/16)\approx1.92379$。
\parta{c,d} 从序列可数出长度，所以 $H(Y|X)=0$。两种链式展开给 $H(X|Y)=H(X)-H(Y)\approx3.88871$。
\problem{4.3}{已知 $\V X=\E[X^2]-(\E X)^2$。用 Jensen 证明方差非负。}
\solve $f(t)=t^2$ 的二阶导数为 $2>0$，故凸。Jensen 给 $f(\E X)\le\E f(X)$，于是 $(\E X)^2\le\E[X^2]$，移项即方差 $\ge0$。严格凸给等号条件为 $X$ 几乎必然恒定。
\problem{4.4}{$X,Y$ 是独立公平 Bernoulli。(a) 算 $I(X;Y)$。(b) 令 $Z=X+Y$，算 $I(X;Y|Z)$。(c) 是否违反数据处理不等式？}
\solve (a) 独立给 $H(X|Y)=H(X)=1$，所以 $I=0$。(b) 四个输入对各占 $1/4$。$Z=0,2$ 时 $X$ 已确定；$Z=1$ 概率 $1/2$，内部两对等可能，故 $H(X|Z)=\tfrac12\cdot1$。从 $X=Z-Y$ 得 $H(X|Y,Z)=0$，故条件互信息 $1/2$。(c) 固定 $Y=0$ 则 $Z=X$，所以给定 $Y$ 后 $X,Z$ 仍相关，$X\to Y\to Z$ 的前提不成立。

\topic{Tut 5 · 基本与概率不等式}
\problem{5.1}{证明 (a) $H(X,Y|Z)\ge H(X|Z)$；(b) $I(X,Y;Z)\ge I(X;Z)$；(c) $H(X,Y,Z)-H(X,Y)\le H(X,Z)-H(X)$，并找取等条件；(d) 证明 $I(X;Z|Y)=I(Z;Y|X)-I(Z;Y)+I(X;Z)$。}
\solve (a) 按条件熵链式展开，左减右为 $H(Y|X,Z)\ge0$；等号要求给定 $(X,Z)$ 后能确定 $Y$。(b) 按互信息链式展开，差为 $I(Y;Z|X)\ge0$；等号要求给定 $X$ 后独立。(c) 左、右分别是 $H(Z|X,Y)$、$H(Z|X)$；右减左就是 $I(Y;Z|X)$，取等同 (b)。(d) $I(X,Y;Z)=I(X;Z)+I(Y;Z|X)=I(Y;Z)+I(X;Z|Y)$；移项再用互信息对称性即得。
\problem{5.2}{比较 (a) $H(5X)$ 与 $H(X)$；(b) $I(g(X);Y)$ 与 $I(X;Y)$；(c) $H(X_0|X_{-1})$ 与 $H(X_0|X_{-1},X_1)$；(d) $H(X,Y)/(H(X)+H(Y))$ 与 $1$。}
\solve (a) 乘、除 5 互逆，故熵相等。(b) $Y\to X\to g(X)$，DPI 给前者 $\le$。(c) 前减后为 $I(X_0;X_1|X_{-1})\ge0$，故前者 $\ge$。(d) 分母减分子是 $I(X;Y)\ge0$，分母为正时比值 $\le1$；两熵都为零时比值无定义。
\problem{5.3}{给定概率质量函数 $p(x)$，证明 $P(p(X)\le d)\log(1/d)\le H(X)$，题面写 $d\ge0$。}
\solve 取 $W=\log(1/p(X))\ge0$，其期望是 $H(X)$。$0<d<1$ 时，$p(X)\le d$ 等价于 $W\ge\log(1/d)>0$。用 Markov 得 $P(p(X)\le d)\le H(X)/\log(1/d)$，乘回阈值。$d\ge1$ 左端非正。$d=0$ 原式无定义；有限熵时，写 $t=\log(1/d)$，有 $tP(W\ge t)\le\E[W1_{\{W\ge t\}}]\to0$，可按极限补为 0。
\form{5.4} 与 7.2 相同，共用下方完整题解。

\topic{Tut 6 · 熵分组、KL、Markov 链}
\problem{6.1}{分布 $p=(p_1,\ldots,p_m)$，把末两项合并成 $q$，新末项为 $s=p_{m-1}+p_m$。证明 $H(p)=H(q)+sH(p_{m-1}/s,p_m/s)$。}
\solve 设 $Y$ 为原结果 $X$ 的组标签。$H(Y|X)=0$，链式得 $H(X)=H(Y)+H(X|Y)$。$H(Y)=H(q)$。未合并的组内只有一个结果，条件熵为 0；合并组用原概率除以组和 $s$，该组熵乘出现概率 $s$，得到右边。若 $s=0$，该组贡献按 0 计。
\problem{6.2}{对联合分布 $p(x,y,z)$，证明 (a) $I(X,Y;Z)-I(Y,Z;X)=I(Y;Z)-I(X;Y)$；(b) $D(p_{XYZ}\Vert p_Xp_Yp_Z)=I(X;Y)+I(X,Y;Z)$。}
\solve (a) 两项依次展开为 $I(Y;Z)+I(X;Z|Y)$ 与 $I(Y;X)+I(Z;X|Y)$。条件项相等，相减抵消。
\parta{b} 在联合概率正的格子，$p(x,y,z)/[p(x)p(y)p(z)]$ 等于 $p(x,y,z)/[p(x,y)p(z)]$ 乘 $p(x,y)/[p(x)p(y)]$。取对数后按联合分布求和，第一项就是 $I(X,Y;Z)$。第二项对数与 $z$ 无关，先算 $\sum_zp(x,y,z)=p(x,y)$，余项就是 $I(X;Y)$。
\problem{6.3}{马尔可夫链 $X_1\to X_2\to X_3$ 的状态数依次为 $n,k,m$，其中 $k<n,k<m$。(a) 求 $I(X_1;X_3)$ 的上界。(b) $k=1$ 时求其值及独立性。}
\solve (a) 数据处理给 $I(X_1;X_3)\le I(X_1;X_2)$。后者等于 $H(X_2)-H(X_2|X_1)\le H(X_2)$；与均匀分布的 KL 给 $\log k-H(X_2)\ge0$，故上界 $\log k$。(b) $k=1$ 时上界为 0；互信息非负，所以等于 0，两端独立。
\problem{6.4}{正面概率 $1/5$ 的硬币独立抛偶数 $N$ 次。用 (a) Markov，(b) Chebyshev 给「至少一半正面」求上界。}
\solve 令 $K\sim\mathrm{Bin}(N,1/5)$，$\E K=N/5$，$\V K=N(1/5)(4/5)=4N/25$。(a) 阈值 $N/2$，故 $P(K\ge N/2)\le(N/5)/(N/2)=2/5$。(b) $K\ge N/2$ 蕴含 $|K-N/5|\ge N/2-N/5=3N/10$，所求概率不超过 $(4N/25)/(3N/10)^2=16/(9N)$；概率界可与 1 取较小值。

\topic{Tut 7 · AEP、计数与定长编码}
\problem{7.1}{正面率 $p$ 的硬币独立抛 $n$ 次。(a) 算正面不少于 $n-1$ 次的精确概率，(b) 用 Markov 求上界，(c) $n=2$ 时上界在精确值的 1\% 以内，求 $p$。}
\solve $K\sim\mathrm{Bin}(n,p)$。(a) 只有 $K=n-1,n$，相加得 $np^{n-1}(1-p)+p^n$。(b) $n\ge2$ 时阈值为正，界为 $\E K/(n-1)=np/(n-1)$，也可与 1 取最小；$n=1$ 时事件必然。
\parta{c} $n=2$ 时精确值 $2p-p^2$，原 Markov 界 $2p$，差 $p^2$。按相对误差，$p>0$ 时 $p^2/(2p-p^2)=p/(2-p)\le.01$，得 $101p\le2$；$p=0$ 两值相同。若指 1 个百分点，$p^2\le.01$，得 $p\le.1$。若先截断上界，在 $p\ge1/2$ 段差为 $(1-p)^2$；相对误差 $(1-p)^2/[1-(1-p)^2]\le.01$ 给 $p\ge1-1/\sqrt{101}$；绝对误差 $(1-p)^2\le.01$ 给 $p\ge.9$。
\problem{7.2 = 5.4}{i.i.d. 二元信源 $p(1)=.005$、$p(0)=.995$，每 100 位为一块。仅给至多 3 个 1 的块编码。(a) 定长码至少几位？(b) 未编码块的概率？(c) 用 Chebyshev 给上界并比较。}
\solve 令 $K$ 为块内 1 的数目。
\parta{a} 恰 $k$ 个 1 只需选位置，有 $\binom{100}k$ 串。编码数为 $1+100+4950+161700=166751$。因为 $2^{17}<166751\le2^{18}$，故最短 18 位。
\parta{b} 每个含 $k$ 个 1 的串概率 $.005^k .995^{100-k}$，乘串数给 $P(K=k)$。未编码是 $K\ge4$，概率 $1-\sum_{k=0}^3\binom{100}k(.005)^k(.995)^{100-k}\approx.00167327$。
\parta{c} $\mu=100(.005)=.5$，$\sigma^2=100(.005)(.995)=.4975$。若 $K\ge4$ 则 $|K-.5|\ge3.5$，故 $P(K\ge4)\le.4975/3.5^2\approx.0406122$；上界大于精确值。
% Keep the full typical-set exercise on the second side of the compact edition.
\ifdefined\LargePrint\else
  \end{multicols}
  \newpage
  \begin{multicols}{\cols}
  \bodyformat\raggedright
\fi
\problem{7.3}{二元信源 $(p(0),p(1))=(.4,.6)$。(a) 求 $H(X)$。(b) 扩展长度 $N=25$，$\beta=.1$。求 (1) 哪些串典型，(2) 典型集合总概率，(3) 典型串数，(4) 最小 $.9$-充分集合大小，(5) 其 essential bit content。}
\solve (a) $H=-.4\log .4-.6\log .6\approx.97095$。
\parta{b1} 若有 $k$ 个 1，单串概率 $.6^k .4^{25-k}$。每符号自信息为 $-(k/25)\log .6-(1-k/25)\log .4$；减去 $H$ 并合并对数，得 $(.6-k/25)\log(1.5)$。要求绝对值小于 $.1$，即 $15-2.5/\log(1.5)<k<15+2.5/\log(1.5)$；边界约 $10.73,19.27$，所以 $k=11,\ldots,19$。
\parta{b2,b3} 每组有 $\binom{25}k$ 串。乘单串概率再加组得 $\sum_{k=11}^{19}\binom{25}k(.6)^k(.4)^{25-k}\approx.936246$；只加各组串数得 $26366510$。
\parta{b4,b5} $k$ 每加 1，单串概率乘 $1.5$；贪心取最高概率串，任何低概率入选串可换成未选的高概率串来增大覆盖率。先取 $k\ge13$ 的 $\sum_{k=13}^{25}\binom{25}k=16777216$ 串，覆盖率 $\sum_{k=13}^{25}\binom{25}k(.6)^k(.4)^{25-k}\approx.84623223$。再取 $k=12$ 的 $\lceil(.9-P(K\ge13))/(.6^{12}.4^{13})\rceil=3680673$ 串，共 $20457889$。于是 $H_{.9}=\log 20457889\approx24.28615$ bits，实际定长编号向上取整为 25 位。

\topic{Tut 8 · 二元 Huffman}
\problem{8.1}{哪组不可能是任何分布的二元 Huffman 码？(a) $\{0,10,11\}$；(b) $\{00,01,10,110\}$；(c) $\{01,10\}$。}
\solve (a) 取概率 $(.5,.3,.2)$，先并后两项为 $.5$，再连根；根左得 $0$，右边两叶得 $10,11$。(b) 有三个长 2 和一个长 3，Kraft 和 $3/4+1/8=7/8$；(c) 两个长 2，和 $1/2$。Huffman 每次合并两项，内部节点都有两子树，因此 Kraft 和为 1，(b,c) 不可能。
\problem{8.2}{对任意 $\epsilon>0$，构造分布使最优平均码长 $L>H(X)+1-\epsilon$。}
\solve 取两符号概率 $(1-\delta,\delta)$，其中 $0<\delta<1/2$。两个码字不能用空串，码 $0,1$ 达到最短长度各 1，故 $L=1$。令 $\delta\downarrow0$，两项 $-\delta\log\delta$ 和 $-(1-\delta)\log(1-\delta)$ 均趋 0，故可取 $\delta$ 使 $H<\epsilon$，代入得 $1>H+1-\epsilon$。
\problem{8.3}{(a) 为 $(.3,.3,.2,.1,.1)$ 构造二元 Huffman 码并求平均长。(b) 找新分布使同一个码的平均码长等于新分布的熵。}
\solve (a) 每轮并最小者：$.1+.1=.2$，$.2+.2=.4$，$.3+.3=.6$，$.4+.6=1$。根左放 $.6$ 的两叶，得 $00,01$；右放 $.4$，得 $10$ 和 $110,111$。平均 $L=2(.3+.3+.2)+3(.1+.1)=2.2$。
\parta{b} 让每符号自信息等于对应长度，即 $p'_i=2^{-\ell_i}$。于是 $p'=(1/4,1/4,1/4,1/8,1/8)$，和为 1，$H=L=3(1/4)2+2(1/8)3=2.25$。
\problem{8.4}{概率 $p_1\ge\dots\ge p_{10}$ 的二元 Huffman 码长为 $\ell_1,\ldots,\ell_{10}$。把末项拆成 $\alpha p_{10},(1-\alpha)p_{10}$，$0\le\alpha\le1$。问新的最优码长。}
\solve 两项都不超过 $p_{10}$，可作为最小对先合并；合回 $p_{10}$ 后，剩下的问题与原树相同。取原来第 10 叶并延展一个分支位给两个孩子，得到一组最优长度 $(\ell_1,\ldots,\ell_9,\ell_{10}+1,\ell_{10}+1)$。新增平均位数为 $\alpha p_{10}+(1-\alpha)p_{10}=p_{10}$。若端点时删掉零概率符号，则回到原码。

\topic{Tut 9 · Shannon 与最优码}
\problem{9.1}{概率 $(1/3,1/3,1/4,1/12)$。(a) 构造 Huffman 码。(b) 证明 $(1,2,3,3)$ 和 $(2,2,2,2)$ 都最优。(c) 用例子说明最优码中的某个长度可大于 Shannon 长度。}
\solve (a,b) 先并 $1/12+1/4=1/3$，余三个同权节点。合一原节点与新节点，取 $0,10,110,111$，长度 $(1,2,3,3)$，平均 $1/3+2/3+3(1/4+1/12)=2$。若改合两原节点，四叶都在第 2 层，平均 2。两种打破并列的方式均执行 Huffman，故都最优。(c) 第三符号 $p=1/4$ 的 Shannon 长是 $\lceil\log4\rceil=2$，第一套却用 3 位。
\problem{9.2}{符号概率为权重 $(1,2,3,4,5,6)$ 分别除以 21。(a) 二元 Huffman，(b) 三元 Huffman，(c) 算两者平均码长。}
\solve (a) 二元并 $1+2=3$、$3+3=6$、$4+5=9$、$6+6=12$、$9+12=21$。依原顺序读叶，例 $1110,1111,110,00,01,10$，长 $(4,4,3,2,2,2)$。
\parta{b} 完整三叉树叶数须奇，补零。并 $0+1+2=3$，再 $3+3+4=10$，最后 $5+6+10=21$。去掉虚叶，一组码是 $201,202,21,22,0,1$。
\parta{c} 二元平均为 $(4+8+9+8+10+12)/21=51/21$ bit/符号；三元为 $(3+6+6+8+5+6)/21=34/21$ 三元位/符号。
\problem{9.3}{求使 $\sum_i p_i\ell_i$ 最小的二元即时码长度。(a) $p=(10,9,8,7,7)/41$；(b) $p_i=.9(.1)^{i-1}$，$i=1,2,\ldots$。}
\solve (a) 合并 $7+7=14$、$8+9=17$、$10+14=24$、$17+24=41$，得长 $(2,2,2,3,3)$，平均 $(2\cdot27+3\cdot14)/41=96/41$。
\parta{b} 取 $0,10,110,1110,\ldots$，长 $i$，平均等于几何分布的均值 $1/.9=10/9$。若首项用至少两位，单项贡献 $\ge1.8>10/9$，可排除。首项用一位后，剩余概率 $.1$ 都在另一分支；扣掉共同首位并归一化，分布又是原几何分布。设最优下确界为 $L_*$，所有符号先花一位，尾部至少贡献 $.1L_*$，故 $L_*\ge1+.1L_*$，得 $L_*\ge10/9$，构造达到。
\problem{9.4}{概率 $p_1\ge\dots\ge p_m>0$，$F_i=\sum_{k<i}p_k$，Shannon 码取 $F_i$ 的前 $\ell_i=\lceil\log(1/p_i)\rceil$ 个二进制小数位。(a) 证明前缀性质与 $H(X)\le L<H(X)+1$。(b) 给 $(1/2,1/4,1/8,1/8)$ 构造码。}
\solve (a) $i<j$ 则 $\ell_i\le\ell_j$，且 $F_j-F_i\ge p_i\ge2^{-\ell_i}$。共享 $\ell_i$ 位前缀等于落在同一宽 $2^{-\ell_i}$ 的半开区间，间距不允许，故前缀性质成立。取整给 $-\log p_i\le\ell_i<-\log p_i+1$，乘 $p_i$ 逐项相加即得长度界。(b) $F=(0,1/2,3/4,7/8)$，二进制为 $.0,.10,.110,.111$；长度 $(1,2,3,3)$，取位得 $0,10,110,111$。

\topic{Tut 10 · 信道容量}
\problem{10.1}{输入输出均为 $\{0,\ldots,10\}$，信道 $Y=X+Z\bmod11$，$Z$ 与 $X$ 独立，在 $1,2,3$ 上均匀。(a) 求容量，(b) 给一个最优输入分布。}
\solve 固定 $X=x$，三个输出是 $x+1,x+2,x+3$ 模 11，各占 $1/3$，故 $H(Y|X)=\log3$。$Y$ 最多 11 种，$H(Y)\le\log11$。取 $X$ 均匀，则 $P(Y=y)=\sum_{z=1}^3(1/3)(1/11)=1/11$，输出熵达上界。故 $C=\log11-\log3=\log(11/3)$，均匀输入最优。
\problem{10.2}{二元输入信道，每次以 $\alpha$ 擦除、以 $\epsilon$ 翻转、以 $1-\alpha-\epsilon$ 正确输出。(a) 求容量，(b) 取 $\alpha=0$，(c) 取 $\epsilon=0$。}
\solve (a) 当 $\alpha<1$，令 $E$ 表示擦除。$Y$ 决定 $E$，而擦除概率不随输入变化，故 $I(X;Y)=I(X;E)+I(X;Y|E)$，首项为 0。擦除时 $Y=e$ 恒定；未擦除概率 $1-\alpha$，翻转条件概率 $q=\epsilon/(1-\alpha)$。该分支的 $H(Y|X,E=0)=\B(q)$，输出熵至多 1，均匀输入做到 1。故 $C=(1-\alpha)[1-\B(q)]$。$\alpha=1$ 时恒擦除，容量 0。(b) $C=1-\B(\epsilon)$。(c) $C=1-\alpha$。
\problem{10.3}{$Y=XZ$，$X,Z$ 是独立二元变量，$P(Z=1)=\alpha$。(a) 求容量及最优输入。(b) 若接收方也能观察 $Z$，再求容量。}
\solve (a) 设 $p=P(X=1)$。$P(Y=1)=\alpha p$，且 $X=0$ 时输出恒零、$X=1$ 时输出是 $Z$，所以 $I=\B(\alpha p)-p\B(\alpha)$。$\alpha=0$ 时容量 0，任意输入均可。其余令 $d=\B(\alpha)/\alpha$。令导数 $\alpha\log[(1-\alpha p)/(\alpha p)]-\B(\alpha)=0$。解得 $t=\alpha p^*=1/(1+2^d)$，$p^*=t/\alpha$。导数随 $p$ 严格下降，近 0 时为正；$0<\alpha<1$ 时，在 $p=1$ 导数为 $\alpha\log[(1-\alpha)/\alpha]-\B(\alpha)=\log(1-\alpha)<0$；$\alpha=1$ 时最优为 $p=1/2$。故 $p^*$ 为最大点。由 $d=\log[(1-t)/t]$，代回 $C=\B(t)-td=-\log(1-t)=\log(1+2^{-d})$。
\parta{b} 链式法则 $I(X;Y,Z)=I(X;Z)+I(X;Y|Z)$。独立使第一项为 0；$Z=0$ 不透露输入，$Z=1$ 时 $Y=X$。故为 $\alpha H(X)$，均匀输入给 $C=\alpha$。
\problem{10.4}{证明给信道转移矩阵增加一列不会降低容量。}
\solve 每列是一个输入的输出分布。给新增输入概率 0，其余输入概率按原最优分布设置。所有旧输入和输出的联合概率照旧，故可达到旧容量。新容量是对扩大后的全部输入分布取最大值，因而 $C_{\rm new}\ge C_{\rm old}$。

\topic{Tut 11 · 容量、错误率与译码}
\form{本组共用的码率推导} 均匀独立消息位 $M_1,\ldots,M_K$，译为 $\hat M$，第 $i$ 位错率 $e_i$，平均 $e$。误差位 $M_i\oplus\hat M_i$ 加上 $\hat M_i$ 就能反推 $M_i$，故 $H(M_i|\hat M_i)\le\B(e_i)$。链式法则和条件降熵给 $H(M|\hat M)\le\sum_iH(M_i|\hat M_i)$，Jensen 给 $\sum_i\B(e_i)\le K\B(e)$。于是 $K[1-\B(e)]\le I(M;\hat M)$。数据处理给 $I(M;\hat M)\le I(X^N;Y^N)$；无记忆给 $H(Y^N|X^N)=\sum_i H(Y_i|X_i)$，而 $H(Y^N)\le\sum_iH(Y_i)$，故 $I(X^N;Y^N)\le NC$。两式合并得 $R\le C/[1-\B(e)]$（$e<1/2$）。这是渐近可逼近的上界：均匀二元信源在 Hamming 失真 $e$ 下的率失真函数为 $1-\B(e)$，即每消息位可压到接近这么多比特，再用低于 $C$ 的信道码传送。
\problem{11.1}{BSC 翻转率 $q$，输入概率 $(p,1-p)$。(a) 算互信息，证明 $C=1-\B(q)$。(b) 允许 5\% 错误概率时，求可达码率的最紧上界。}
\solve (a) 输出 0 来自输入 0 未翻转或输入 1 翻转，故 $p_Y(0)=p(1-q)+(1-p)q$。条件输出熵总为 $\B(q)$，所以 $I=\B(q+p(1-2q))-\B(q)$。首项至多 1，$p=1/2$ 时达到，故 $C=1-\B(q)$。(b) 按平均消息位错率解，共用推导给 $R\le[1-\B(q)]/[1-\B(.05)]$，分母约 $.713603$。若题中 5\% 指固定整块错误率，渐近上界由强逆定理给 $C$；若指信道翻转率则代入 $q=.05$，容量约 $.713603$。
\problem{11.2}{BEC 擦除率 $q$，输入概率 $(p,1-p)$。(a) 算互信息并证 $C=1-q$。(b,c) 容许错误概率 25\%、50\% 时分别求最紧码率上界。}
\solve (a) 收到 0 或 1 时输入确定；收到 $e$ 时，由于擦除与输入独立，后验不变。因此 $H(X|Y)=q\B(p)$，互信息 $(1-q)\B(p)$ 在 $p=1/2$ 取 $C=1-q$。(b) 平均消息位错率 $.25$ 给 $R\le(1-q)/[1-\B(.25)]\approx(1-q)/.188722$。(c) 允许 $.5$ 时，什么都不传，对均匀位盲猜的平均错误率就为 $.5$，故此条件不给有限码率上界。若百分比指固定整块错误率，两问的渐近上界都是 $C$；若指信道擦除率，容量分别是 $.75$、$.5$。
\problem{11.3}{输入 $a,b,c$，输出 $a,b,c,d$，转移矩阵四行是 $(3,1,0)/4$、$(1,2,1)/4$、$(0,1,2)/4$、$(0,0,1)/4$。输入概率 $(1/4,1/4,1/2)$。(a) 算互信息。(b) 对每个输出用最大 $P(y|x)$ 的输入译码，算平均错误率。(c,d) 码率分别为 $.01$ 和 $100$ bit/次、位误率 1\% 的声明可能吗？}
\solve (a) 矩阵乘输入：$p_Y(a)=3/16+1/16=1/4$；$p_Y(b)=1/16+2/16+1/8=5/16$；$p_Y(c)=1/16+1/4=5/16$；$p_Y(d)=1/8$。第一列非零概率为 $(3/4,1/4)$，后两列均有 $(1/2,1/4,1/4)$。故 $H(Y|X)=\tfrac14\B(1/4)+\tfrac34(1.5)$；$H(Y)=\tfrac14(2)-2(\tfrac5{16})\log(\tfrac5{16})+\tfrac18(3)$。相减得 $I\approx.59598$。
\parta{b} 按每行最大值，$a,b,c,d$ 分别译成 $a,b,c,c$。每种正确事件用先验乘该条件概率，正确率为 $3/16+1/8+1/4+1/8=11/16$，错率 $5/16$。
\parta{c,d} (a) 的输入已达到 $I>.01$，故 $C>.01$，足够长的码可使整块、从而平均位错误率都低于 1\%。另一方面 $C\le H(X)\le\log3$ 对最大化取值仍成立。共用推导给 $R\le\log3/[1-\B(.01)]\approx1.7243$；100 不可能。
\problem{11.4}{独立加性噪声信道 $Y=X+Z$，$X\in\{0,1\}$，噪声等概率取 $0,a$，$a$ 为任意整数。求容量。}
\solve 给定输入 0、1 时，输出集合分别为 $\{0,a\},\{1,1+a\}$。交集非空仅当 $a=1$（共享 1）或 $a=-1$（共享 0）。共享输出在两输入下都以 $1/2$ 出现，后验保持先验；其它输出确定输入。于是 $H(X|Y)=\tfrac12H(X)$，互信息最大为 $1/2$。其它 $a$ 时集合不交（$a=0$ 按噪声恒 0 理解），所以 $H(X|Y)=0$，均匀输入给 $C=1$。

\topic{Tut 12 · Hamming 与算术编码}
\problem{12.1}{对 $(7,4)$ Hamming 码，(a) 编码 \bb{0111}；(b) 用题目所指维恩图解码 \bb{1100000}；(c) 给该串的校验综合；(d) 最多保证纠正几位？同样 4 个数据位，要用多少位的重复码才能保证相同能力？}
\solve 原题 PDF 缺图，采用常见偶校验约定：数据位置 $3,5,6,7$，三组位置 $(1,3,5,7)$、$(2,3,6,7)$、$(4,5,6,7)$。
\parta{a} 数据位依次 $0,1,1,1$。第一组其余位异或 $0\oplus1\oplus1=0$，第二组 $0\oplus1\oplus1=0$，第三组 $1\oplus1\oplus1=1$，故校验位 $(p_1,p_2,p_4)=(0,0,1)$，组装得 \bb{0001111}。
\parta{b,c} 对 \bb{1100000} 按三组求实际奇偶，得到 $(s_1,s_2,s_4)=(1,1,0)$。综合按高位在前为 \bb{011}，即错误位置 $1+2=3$。翻第 3 位得 \bb{1110000}，取数据位得 \bb{1000}。
\parta{d} 校验矩阵的七列是不同非零的三位位置编号，单列或两列异或都不为零，所以非零码字至少三位；编号 $1,2,3$ 的列异或为零，距离可达 3。半径 1 的译码球互斥，半径 2 的可能重叠，保证纠一位。重复两次遇一错会平票，三次多数决可纠一错，发 $4\cdot3=12$ 位。
\problem{12.2}{首位 0、1 各半。之后下一个符号与前一位相同的概率 $3/5$，不同的概率 $1/5$，结束 EOS 概率 $1/5$。(a) 算术编码 \bb{011}；(b) 解回；(c) 自选一串编码并解回。}
\solve 按 $0,1,\code$ 排序，首轮只有 $0,1$。
\parta{a} 首位 $0$ 取 $[0,.5)$。前位为 0 时，$1$ 的累积段是 $[.6,.8)$，映到 $[.3,.4)$。前位成了 1，$1$ 的累积段是 $[.2,.8)$，映到 $[.32,.38)$。再取 EOS 的最后 20\%，区间 $[.32+.06(.8),.38)=[.368,.38)$。$[95/256,96/256)$ 在其中，所以发 95 的 8 位二进制 \bb{01011111}。
\parta{b} 代表值 $r=95/256=.37109375$。全区间 $r<.5$，读 0；在 $[0,.5)$ 中，$r$ 落在 1 段 $[.3,.4)$；该区间中落在 1 段 $[.32,.38)$；再落在 EOS 段 $[.368,.38)$，读出 \bb{011}。
\parta{c} 选 \bb{0}。首符号取 $[0,.5)$，EOS 切出 $[.4,.5)$。取 $[7/16,8/16)$，发 \bb{0111}。$r=7/16$ 先落 0 段，再落 $[.4,.5)$ 的 EOS 段，译为 \bb{0}。
\problem{12.3}{(a) 为何使用自适应分布？(b) Dirichlet 分布的直观含义？(c) 符号 $a,b,c$ 的虚拟计数为 $(3,1,1)$，EOS 固定 $1/4$。以自适应算术码编码 \bb{bb}。(d) 解回。}
\solve (a) 源的概率未知或会变化时，按观察更新，常见符号获得较宽的区间；编解码同步使用同一更新规则。(b) Dirichlet 是对概率向量的分布。参数可看作虚拟计数；每观察某符号，其计数加 1，预测概率是该计数占总计数的比例。
\parta{c} 选顺序 $a,b,c,\code$。前三项的计数比乘 $3/4$。开始 $a,b,c$ 的概率为 $.45,.15,.15$，故第一个 $b$ 取 $[.45,.60)$。更新 $(3,2,1)$ 后，$a,b$ 概率 $.375,.25$；第二个 $b$ 取 $[.45+.15(.375),.45+.15(.625))=[.50625,.54375)$。更新 $(3,3,1)$；EOS 仍是末 $1/4$，起点 $.50625+.0375(.75)=.534375$，终点 $.54375$。选 $[137/256,138/256)$，发 \bb{10001001}。
\parta{d} $r=137/256=.53515625$。首轮在 $[.45,.60)$，译 $b$ 并更新 $(3,2,1)$；第二轮在 $[.50625,.54375)$，再译 $b$，更新 $(3,3,1)$；第三轮在 $[.534375,.54375)$，读 EOS 停止，输出 \bb{bb}。
\end{multicols}
\end{document}
'
\documentclass[a4paper]{article}
\usepackage[UTF8,scheme=plain,fontset=fandol]{ctex}
\usepackage{amsmath,amssymb,multicol,xcolor}
\ifdefined\LargePrint
  \usepackage[margin=14mm]{geometry}
  \setlength{\columnsep}{6mm}
  \newcommand{\bodyformat}{\fontsize{9.6}{13}\selectfont}
  \newcommand{\cols}{2}
  \setlength{\parskip}{2.7pt}
\else
  \usepackage[top=5.6mm,bottom=7.6mm,left=5.8mm,right=5.8mm]{geometry}
  \setlength{\columnsep}{3mm}
  \newcommand{\bodyformat}{\fontsize{6.65}{7.8}\selectfont}
  \newcommand{\cols}{4}
  \setlength{\parskip}{1.05pt}
\fi
\usepackage{fancyhdr}
\definecolor{ink}{HTML}{163D51}
\pagestyle{fancy}
\fancyhf{}
\fancyfoot[L]{\scriptsize COMP2610 / COMP6261}
\fancyfoot[R]{\scriptsize \thepage}
\renewcommand{\headrulewidth}{0pt}
\setlength{\footskip}{4.5mm}
\setlength{\parindent}{0pt}
\setlength{\multicolsep}{2pt}
\setlength{\emergencystretch}{1.3em}
\setlength{\columnseprule}{.14pt}
\setlength{\premulticols}{0pt}
\setlength{\postmulticols}{0pt}
\raggedcolumns
\newcommand{\topic}[1]{\par\addvspace{2.6pt}{\color{ink}\bfseries #1}\par\nobreak}
\newcommand{\problem}[2]{\par\addvspace{1.4pt}\textcolor{ink}{\textbf{#1}}\enspace\textbf{题}\enspace #2\par\nobreak}
\newcommand{\solve}{\noindent\textbf{解}\enspace}
\newcommand{\parta}[1]{\par\noindent\textbf{(#1)}\enspace}
\newcommand{\form}[1]{\par\noindent\textbf{#1}\enspace}
\newcommand{\E}{\mathbb E}
\newcommand{\V}{\operatorname{Var}}
\newcommand{\B}{h_2}
\newcommand{\bb}[1]{\texttt{#1}}
\newcommand{\code}{\mathrm{EOS}}
\begin{document}
{\color{ink}\fontsize{13}{15}\selectfont\bfseries COMP2610 / COMP6261 \quad L22 公式与 Tut 逐题解}\par
{\fontsize{7}{8.4}\selectfont Tut 1–12 共 46 道去重题。题干以中文保留条件和小问；Tut 5 Q4 与 Tut 7 Q2 共用一份解答。信息量对数以 2 为底。\par}
\begin{multicols}{\cols}
\bodyformat\raggedright
\topic{公式区}
\form{概率与 Bayes} $p(x,y)=p(x)p(y|x)$，$p(y)=\sum_xp(x,y)$，$p(x|y)=p(x,y)/p(y)$。独立时 $p(x,y)=p(x)p(y)$；条件独立在固定条件后用同一乘积形式。
\form{Expected Value} $\E g(X)=\sum_x p(x)g(x)$。期望可直接相加。
\form{Variance} $\V X=\E[X^2]-(\E X)^2$。独立变量相加时方差相加。
\form{Binomial Distribution} $K=\sum_{i=1}^n B_i$，$B_i$ 独立且为 $\mathrm{Bern}(p)$。$P(K=k)=\binom nk p^k(1-p)^{n-k}$，$\E K=np$，$\V K=np(1-p)$。理由：$B_i^2=B_i$，所以单次期望 $p$、方差 $p-p^2$。
\form{Markov's Inequality} $W\ge0,\ a>0$ 时，$P(W\ge a)\le\E W/a$。阈值以上每个结果至少贡献 $a$。
\form{Chebyshev's Inequality} $\mu=\E W,\ t>0$ 时，$P(|W-\mu|\ge t)\le\V(W)/t^2$。单侧阈值 $a>\mu$ 用 $t=a-\mu$。这是对 $(W-\mu)^2$ 用 Markov。
\form{Jensen / MLE} 凸 $f$ 满足 $f(\E X)\le\E f(X)$。i.i.d. 的 MLE 最大化 $\ln L(\theta)=\sum_i\ln p_\theta(x_i)$。
\form{Entropy} 自信息 $-\log p(x)$；$H(X)=-\sum_xp(x)\log p(x)$；$H(X,Y)=-\sum_{x,y}p(x,y)\log p(x,y)$。$H(Y|X)=\sum_xp(x)H(Y|X=x)$。$0\log0=0$。$H(X,Y)=H(X)+H(Y|X)$；$H(X)\le\log|\mathcal X|$。记 $\B(p)=-p\log p-(1-p)\log(1-p)$。
\form{互信息与 KL} $I(X;Y)=H(X)-H(X|Y)=H(Y)-H(Y|X)=D(p_{XY}\Vert p_Xp_Y)$；$D(p\Vert q)=\sum p\log(p/q)\ge0$。条件互信息是给定条件时互信息的加权平均。对 $X\to Y\to Z$，$I(X;Z)\le\min\{I(X;Y),I(Y;Z)\}$。
\form{Extended Ensemble} $X^N$ 把整块视作一个结果。i.i.d. 时 $p(x^N)=\prod_i p(x_i)$，$H(X^N)=NH(X)$。含 $k$ 个 1 的单串概率为 $p^k(1-p)^{N-k}$，同组有 $\binom Nk$ 串。
\form{AEP／典型集} $T_{N,\beta}=\{x^N:|-\log p(x^N)/N-H(X)|<\beta\}$，i.i.d. 时其概率趋 1。最小 $\delta$-充分集按单串概率降序取至累计概率 $\ge\delta$；$H_\delta=\log|S_\delta|$，定长需 $\lceil H_\delta\rceil$ 位。固定 $0<\delta<1$ 时 $H_\delta/N\to H(X)$。
\form{信源编码} $L=\sum_i p_i\ell_i$；二元 Kraft 和 $\sum_i2^{-\ell_i}\le1$。Huffman 每轮并两个最小权重。Shannon 取累计 $F_i=\sum_{j<i}p_j$ 的前 $\ell_i=\lceil-\log p_i\rceil$ 位；SFE 取 $F_i+p_i/2$ 的前 $\ell_i+1$ 位。$H\le L_{\rm opt}\le L_{\rm Shannon}<H+1$。
\form{算术编码} 当前 $[a,a+w)$，符号概率段 $[F,F+p)$，更新为 $[a+wF,a+w(F+p))$。含 EOS 的最终区间内取一个完整二进制网格区间。
\form{信道编码} $Q_{y,x}=p(y|x)$，$p_Y=Qp_X$；$H(Y|X)=\sum_xp(x)H(Q_{\cdot,x})$；$C=\max_{p_X}I(X;Y)$。$R=K/N$；$P_B=P(\hat{\mathbf s}\ne\mathbf s)$；$R<C$ 可使块错率趋零，$R>C$ 不可。

\topic{Tut 1 · 基础概率}
\problem{1.1}{等分转盘标 $1,2,3,4,5$，求 (a) 转到 4，(b) 转到偶数，(c) 转到质数的概率。}
\solve 五格各占 $1/5$，互斥格可加。(a) $\{4\}$：$1/5$。(b) $\{2,4\}$：$2/5$。(c) $\{2,3,5\}$：$3/5$。
\problem{1.2}{ACT 车牌是三字母接三数字。求随机车牌以 ACT 开头且以 7 结尾的概率。}
\solve 按各合法组合等可能且允许重复，总数 $26^3\cdot10^3$。固定前三字母和末位后，中间两位各有 10 种，共 $10^2$。概率 $10^2/(26^3\cdot10^3)=1/175760$。
\problem{1.3}{四条路的测速启用率 $(.4,.3,.2,.3)$，Arya 选路概率 $(.2,.1,.5,.2)$，且总超速。求 (a) 一早被罚，(b) 五早都不被罚。}
\solve (a) 对每条路先乘选路概率和该路处罚率，再相加，得 $.2(.4)+.1(.3)+.5(.2)+.2(.3)=.27$。(b) 每早不被罚概率 $1-.27=.73$。按各早独立且规则不变，五早全不罚为 $.73^5\approx.20731$。
\problem{1.4}{袋内 5 蓝、3 红、2 黄，抽三球。求红球少于 2 颗的概率，分别按 (a) 放回，(b) 不放回。}
\solve (a) 红概率 $.3$，非红 $.7$，三次独立。零红 $.7^3$；一红的出现位置有 3 个，概率 $3(.3)(.7)^2$；和为 $.784$。(b) 把无序三球子集当等可能结果。零红选三非红，有 $\binom73=35$ 种；一红选一红两非红，有 $\binom31\binom72=63$ 种。总数 $\binom{10}3=120$，故 $98/120=49/60$。
\problem{1.5}{两台录像机独立成功，概率分别为 $.7,.6$。求 (a) 零份，(b) 一份，(c) 两份录像的概率。}
\solve 失败率分别是 $.3,.4$。(a) 两失败：$.3(.4)=.12$。(b) 仅第一台或仅第二台，两情形互斥：$.7(.4)+.3(.6)=.46$。(c) 两成功：$.7(.6)=.42$。三项和为 1。

\topic{Tut 2 · 条件概率与 Bayes}
\problem{2.1}{红、蓝、绿箱被选概率为 $(.2,.2,.6)$，各箱苹果、橙子、青柠数分别为 $(3,4,3)$、$(1,1,0)$、$(3,3,4)$。选箱后均匀抽一果。求 (a) 苹果概率，(b) 给定橙子，来自绿箱的概率。}
\solve (a) 各箱苹果比例 $3/10,1/2,3/10$；乘箱概率后总和为 $.06+.10+.18=.34$。(b) 橙的三项联合概率为 $.2(4/10)=.08$、$.2(1/2)=.10$、$.6(3/10)=.18$。用绿箱项除总和，得 $.18/.36=1/2$。
\problem{2.2}{共 1000 次，行 $Y=0,1$、列 $X=0,1$ 的计数为 $(100,250)$、$(150,500)$。求 (a) $p(1,1)$，(b) $P(X=1)$，(c) $\E X$，(d) $P(Y=1|X=1)$，(e) $P(Y=1|X=0)$。(f) 若 $X\ne Y$ 时 $P(Z=1|X,Y)=.9$，相同则为 $.1$，求 $P(X=1,Y=1|Z=1)$。}
\solve (a) 单格 $500/1000=.5$。(b) $X=1$ 列和 $(250+500)/1000=.75$。(c) $0(.25)+1(.75)=.75$。(d) 固定 $X=1$，分母 750，得 $500/750=2/3$。(e) 固定 $X=0$，分母 250，得 $150/250=.6$。(f) 不同格合计 $400/1000=.4$。全概率 $P(Z=1)=.9(.4)+.1(.6)=.42$。目标格且 $Z=1$ 的联合概率是 $.5(.1)=.05$，所以后验 $=.05/.42=5/42$。
\problem{2.3}{独立抛公平硬币两次，每次正面放白球，反面放红球。随后从这两球中放回抽三次，均为红。求箱内两球都红的后验概率。}
\solve 令红球数 $K$。由两次独立抛币，$K=0,1,2$ 的先验为 $(1,2,1)/4$。固定 $K=k$ 后，放回抽每次红率 $k/2$，三红似然为 $(k/2)^3$。乘先验得未归一化后验 $(0,1/16,1/4)$。故 $P(K=2|\text{三红})=(1/4)/(1/16+1/4)=4/5$。
\problem{2.4}{(a) 从条件概率的定义证明 $p(x|y,z)=p(y|x,z)p(x|z)/p(y|z)$。(b) 证明给定 $Y$ 时仍可对任一变量 $X_i$ 求和边缘化。}
\solve (a) 在分母为正时，$p(x|y,z)=p(x,y|z)/p(y|z)$。分子按乘法规则为 $p(x|z)p(y|x,z)$，代回即得。(b) 令 $U$ 代表除 $X_i$ 外的变量。$p(u|y)=p(u,y)/p(y)=[\sum_x p(u,x,y)]/p(y)=\sum_x p(u,x|y)$。

\topic{Tut 3 · MLE 与熵}
\problem{3.1}{取值 $1,2,3$ 的概率为 $(\theta/2,\theta/2,1-\theta)$，$0\le\theta\le1$。i.i.d. 样本中各频数是 $n_1,n_2,n_3$。(a) 写似然。(b) 样本 $3,3,1,2,3,2,2,1,3,1$ 求 MLE。}
\solve (a) 把单次概率相乘，同类因子合并，得 $L=(\theta/2)^{n_1+n_2}(1-\theta)^{n_3}$。(b) 频数 $(3,3,4)$；$\ln L=6\ln(\theta/2)+4\ln(1-\theta)$。令导数 $6/\theta-4/(1-\theta)=0$，得 $6(1-\theta)=4\theta$，即 $\theta=.6$。导数在此点左右分别正负，两端似然为 0，故为最大值。
\problem{3.2}{联合表的行是 $Y=1,2,3,4$，列是 $X=1,2,3,4$，矩阵按 $1/16$ 为单位，四行为 $(0,0,2,2)$、$(2,1,1,0)$、$(2,2,0,0)$、$(0,1,1,2)$。(a) 证不独立。(b) 求 $H(X)$、$H(Y)$、$H(X|Y)$、$H(Y|X)$、$H(X,Y)$。}
\solve (a) 各行各列都合计 $4/16=1/4$。左上为 0，而独立时应为 $1/4\cdot1/4=1/16$。(b) 两个边缘均匀，两个熵都是 $4(1/4)\log 4=2$。每行除以 $1/4$，非零条件概率为 $(1/2,1/2)$ 或 $(1/2,1/4,1/4)$，熵分别为 $1$ 或 $1/2+1/2+1/2=1.5$。四行熵 $(1,1.5,1,1.5)$，均值 $1.25=H(X|Y)$。四列熵 $(1,1.5,1.5,1)$，均值 $1.25=H(Y|X)$。联合表有六个 $1/8$ 和四个 $1/16$，所以 $H(X,Y)=6(1/8)3+4(1/16)4=3.25$。
\problem{3.3}{均匀抽 52 张牌之一。$V,S$ 为点数和花色，人头牌是 J、Q、K。求 (a) 看到红 K 的自信息，(b) 已知是人头牌时看到 K 的条件自信息，(c) $H(S)$ 与 $H(V,S)$。}
\solve (a) 红 K 有两张，$-\log(2/52)=\log26$。(b) 12 张人头牌中四张 K，$-\log(4/12)=\log3$。(c) 花色各有 $13/52=1/4$，$H(S)=2$；一对 $(V,S)$ 确定一张牌，52 对等可能，$H(V,S)=\log52$。
\problem{3.4}{有限离散 $X$。比较 $H(X)$ 与 $H(Y)$，分别取 (a) $Y=2^X$，(b) $Y=\cos X$。}
\solve 确定映射给 $H(Y|X)=0$，两种顺序展开联合熵得 $H(X)=H(Y)+H(X|Y)$。(a) $X=\log_2Y$，故 $H(X|Y)=0$，两熵相等。(b) $H(X|Y)\ge0$，故 $H(Y)\le H(X)$。若 $X$ 等概率为 $\pm1$，余弦相同，熵从 1 变 0；若在支持集上可逆则取等。

\topic{Tut 4 · 熵、Jensen、互信息}
\problem{4.1}{几何变量 $Y$ 的概率为 $p(1-p)^{y-1}$，$y=1,2,\dots$。求期望和方差。}
\solve 设 $q=1-p$、$m=\E Y$、$s=\E Y^2$。首轮成功时 $Y=1$；失败时 $Y=1+Y'$，$Y'$ 同分布。故 $m=p+q(1+m)=1+qm$，得 $m=1/p$。平方后 $s=p+q(1+2m+s)=1+2qm+qs$，故 $s=(1+2q/p)/p=(2-p)/p^2$；方差 $s-m^2=(1-p)/p^2$。
\problem{4.2}{A、B 两队独立且各半胜率，先赢四局者获胜。$X$ 是完整胜负序列，$Y$ 是总局数。求 (a) $H(X)$，(b) $H(Y)$，(c) $H(Y|X)$，(d) $H(X|Y)$。}
\solve 选总冠军有两种；总长 $y$ 时第 $y$ 局必由冠军赢，前 $y-1$ 局选出其前三胜，故 $M_y=2\binom{y-1}3$。$y=4,5,6,7$ 的序列数为 $2,8,20,40$。每序列概率 $2^{-y}$，所以 $p_Y=(1/8,1/4,5/16,5/16)$。
\parta{a} 序列自信息为 $y$，故 $H(X)=4/8+5/4+6(5/16)+7(5/16)=5.8125$。
\parta{b} $H(Y)=(1/8)3+(1/4)2-2(5/16)\log(5/16)\approx1.92379$。
\parta{c,d} 从序列可数出长度，所以 $H(Y|X)=0$。两种链式展开给 $H(X|Y)=H(X)-H(Y)\approx3.88871$。
\problem{4.3}{已知 $\V X=\E[X^2]-(\E X)^2$。用 Jensen 证明方差非负。}
\solve $f(t)=t^2$ 的二阶导数为 $2>0$，故凸。Jensen 给 $f(\E X)\le\E f(X)$，于是 $(\E X)^2\le\E[X^2]$，移项即方差 $\ge0$。严格凸给等号条件为 $X$ 几乎必然恒定。
\problem{4.4}{$X,Y$ 是独立公平 Bernoulli。(a) 算 $I(X;Y)$。(b) 令 $Z=X+Y$，算 $I(X;Y|Z)$。(c) 是否违反数据处理不等式？}
\solve (a) 独立给 $H(X|Y)=H(X)=1$，所以 $I=0$。(b) 四个输入对各占 $1/4$。$Z=0,2$ 时 $X$ 已确定；$Z=1$ 概率 $1/2$，内部两对等可能，故 $H(X|Z)=\tfrac12\cdot1$。从 $X=Z-Y$ 得 $H(X|Y,Z)=0$，故条件互信息 $1/2$。(c) 固定 $Y=0$ 则 $Z=X$，所以给定 $Y$ 后 $X,Z$ 仍相关，$X\to Y\to Z$ 的前提不成立。

\topic{Tut 5 · 基本与概率不等式}
\problem{5.1}{证明 (a) $H(X,Y|Z)\ge H(X|Z)$；(b) $I(X,Y;Z)\ge I(X;Z)$；(c) $H(X,Y,Z)-H(X,Y)\le H(X,Z)-H(X)$，并找取等条件；(d) 证明 $I(X;Z|Y)=I(Z;Y|X)-I(Z;Y)+I(X;Z)$。}
\solve (a) 按条件熵链式展开，左减右为 $H(Y|X,Z)\ge0$；等号要求给定 $(X,Z)$ 后能确定 $Y$。(b) 按互信息链式展开，差为 $I(Y;Z|X)\ge0$；等号要求给定 $X$ 后独立。(c) 左、右分别是 $H(Z|X,Y)$、$H(Z|X)$；右减左就是 $I(Y;Z|X)$，取等同 (b)。(d) $I(X,Y;Z)=I(X;Z)+I(Y;Z|X)=I(Y;Z)+I(X;Z|Y)$；移项再用互信息对称性即得。
\problem{5.2}{比较 (a) $H(5X)$ 与 $H(X)$；(b) $I(g(X);Y)$ 与 $I(X;Y)$；(c) $H(X_0|X_{-1})$ 与 $H(X_0|X_{-1},X_1)$；(d) $H(X,Y)/(H(X)+H(Y))$ 与 $1$。}
\solve (a) 乘、除 5 互逆，故熵相等。(b) $Y\to X\to g(X)$，DPI 给前者 $\le$。(c) 前减后为 $I(X_0;X_1|X_{-1})\ge0$，故前者 $\ge$。(d) 分母减分子是 $I(X;Y)\ge0$，分母为正时比值 $\le1$；两熵都为零时比值无定义。
\problem{5.3}{给定概率质量函数 $p(x)$，证明 $P(p(X)\le d)\log(1/d)\le H(X)$，题面写 $d\ge0$。}
\solve 取 $W=\log(1/p(X))\ge0$，其期望是 $H(X)$。$0<d<1$ 时，$p(X)\le d$ 等价于 $W\ge\log(1/d)>0$。用 Markov 得 $P(p(X)\le d)\le H(X)/\log(1/d)$，乘回阈值。$d\ge1$ 左端非正。$d=0$ 原式无定义；有限熵时，写 $t=\log(1/d)$，有 $tP(W\ge t)\le\E[W1_{\{W\ge t\}}]\to0$，可按极限补为 0。
\form{5.4} 与 7.2 相同，共用下方完整题解。

\topic{Tut 6 · 熵分组、KL、Markov 链}
\problem{6.1}{分布 $p=(p_1,\ldots,p_m)$，把末两项合并成 $q$，新末项为 $s=p_{m-1}+p_m$。证明 $H(p)=H(q)+sH(p_{m-1}/s,p_m/s)$。}
\solve 设 $Y$ 为原结果 $X$ 的组标签。$H(Y|X)=0$，链式得 $H(X)=H(Y)+H(X|Y)$。$H(Y)=H(q)$。未合并的组内只有一个结果，条件熵为 0；合并组用原概率除以组和 $s$，该组熵乘出现概率 $s$，得到右边。若 $s=0$，该组贡献按 0 计。
\problem{6.2}{对联合分布 $p(x,y,z)$，证明 (a) $I(X,Y;Z)-I(Y,Z;X)=I(Y;Z)-I(X;Y)$；(b) $D(p_{XYZ}\Vert p_Xp_Yp_Z)=I(X;Y)+I(X,Y;Z)$。}
\solve (a) 两项依次展开为 $I(Y;Z)+I(X;Z|Y)$ 与 $I(Y;X)+I(Z;X|Y)$。条件项相等，相减抵消。
\parta{b} 在联合概率正的格子，$p(x,y,z)/[p(x)p(y)p(z)]$ 等于 $p(x,y,z)/[p(x,y)p(z)]$ 乘 $p(x,y)/[p(x)p(y)]$。取对数后按联合分布求和，第一项就是 $I(X,Y;Z)$。第二项对数与 $z$ 无关，先算 $\sum_zp(x,y,z)=p(x,y)$，余项就是 $I(X;Y)$。
\problem{6.3}{马尔可夫链 $X_1\to X_2\to X_3$ 的状态数依次为 $n,k,m$，其中 $k<n,k<m$。(a) 求 $I(X_1;X_3)$ 的上界。(b) $k=1$ 时求其值及独立性。}
\solve (a) 数据处理给 $I(X_1;X_3)\le I(X_1;X_2)$。后者等于 $H(X_2)-H(X_2|X_1)\le H(X_2)$；与均匀分布的 KL 给 $\log k-H(X_2)\ge0$，故上界 $\log k$。(b) $k=1$ 时上界为 0；互信息非负，所以等于 0，两端独立。
\problem{6.4}{正面概率 $1/5$ 的硬币独立抛偶数 $N$ 次。用 (a) Markov，(b) Chebyshev 给「至少一半正面」求上界。}
\solve 令 $K\sim\mathrm{Bin}(N,1/5)$，$\E K=N/5$，$\V K=N(1/5)(4/5)=4N/25$。(a) 阈值 $N/2$，故 $P(K\ge N/2)\le(N/5)/(N/2)=2/5$。(b) $K\ge N/2$ 蕴含 $|K-N/5|\ge N/2-N/5=3N/10$，所求概率不超过 $(4N/25)/(3N/10)^2=16/(9N)$；概率界可与 1 取较小值。

\topic{Tut 7 · AEP、计数与定长编码}
\problem{7.1}{正面率 $p$ 的硬币独立抛 $n$ 次。(a) 算正面不少于 $n-1$ 次的精确概率，(b) 用 Markov 求上界，(c) $n=2$ 时上界在精确值的 1\% 以内，求 $p$。}
\solve $K\sim\mathrm{Bin}(n,p)$。(a) 只有 $K=n-1,n$，相加得 $np^{n-1}(1-p)+p^n$。(b) $n\ge2$ 时阈值为正，界为 $\E K/(n-1)=np/(n-1)$，也可与 1 取最小；$n=1$ 时事件必然。
\parta{c} $n=2$ 时精确值 $2p-p^2$，原 Markov 界 $2p$，差 $p^2$。按相对误差，$p>0$ 时 $p^2/(2p-p^2)=p/(2-p)\le.01$，得 $101p\le2$；$p=0$ 两值相同。若指 1 个百分点，$p^2\le.01$，得 $p\le.1$。若先截断上界，在 $p\ge1/2$ 段差为 $(1-p)^2$；相对误差 $(1-p)^2/[1-(1-p)^2]\le.01$ 给 $p\ge1-1/\sqrt{101}$；绝对误差 $(1-p)^2\le.01$ 给 $p\ge.9$。
\problem{7.2 = 5.4}{i.i.d. 二元信源 $p(1)=.005$、$p(0)=.995$，每 100 位为一块。仅给至多 3 个 1 的块编码。(a) 定长码至少几位？(b) 未编码块的概率？(c) 用 Chebyshev 给上界并比较。}
\solve 令 $K$ 为块内 1 的数目。
\parta{a} 恰 $k$ 个 1 只需选位置，有 $\binom{100}k$ 串。编码数为 $1+100+4950+161700=166751$。因为 $2^{17}<166751\le2^{18}$，故最短 18 位。
\parta{b} 每个含 $k$ 个 1 的串概率 $.005^k .995^{100-k}$，乘串数给 $P(K=k)$。未编码是 $K\ge4$，概率 $1-\sum_{k=0}^3\binom{100}k(.005)^k(.995)^{100-k}\approx.00167327$。
\parta{c} $\mu=100(.005)=.5$，$\sigma^2=100(.005)(.995)=.4975$。若 $K\ge4$ 则 $|K-.5|\ge3.5$，故 $P(K\ge4)\le.4975/3.5^2\approx.0406122$；上界大于精确值。
% Keep the full typical-set exercise on the second side of the compact edition.
\ifdefined\LargePrint\else
  \end{multicols}
  \newpage
  \begin{multicols}{\cols}
  \bodyformat\raggedright
\fi
\problem{7.3}{二元信源 $(p(0),p(1))=(.4,.6)$。(a) 求 $H(X)$。(b) 扩展长度 $N=25$，$\beta=.1$。求 (1) 哪些串典型，(2) 典型集合总概率，(3) 典型串数，(4) 最小 $.9$-充分集合大小，(5) 其 essential bit content。}
\solve (a) $H=-.4\log .4-.6\log .6\approx.97095$。
\parta{b1} 若有 $k$ 个 1，单串概率 $.6^k .4^{25-k}$。每符号自信息为 $-(k/25)\log .6-(1-k/25)\log .4$；减去 $H$ 并合并对数，得 $(.6-k/25)\log(1.5)$。要求绝对值小于 $.1$，即 $15-2.5/\log(1.5)<k<15+2.5/\log(1.5)$；边界约 $10.73,19.27$，所以 $k=11,\ldots,19$。
\parta{b2,b3} 每组有 $\binom{25}k$ 串。乘单串概率再加组得 $\sum_{k=11}^{19}\binom{25}k(.6)^k(.4)^{25-k}\approx.936246$；只加各组串数得 $26366510$。
\parta{b4,b5} $k$ 每加 1，单串概率乘 $1.5$；贪心取最高概率串，任何低概率入选串可换成未选的高概率串来增大覆盖率。先取 $k\ge13$ 的 $\sum_{k=13}^{25}\binom{25}k=16777216$ 串，覆盖率 $\sum_{k=13}^{25}\binom{25}k(.6)^k(.4)^{25-k}\approx.84623223$。再取 $k=12$ 的 $\lceil(.9-P(K\ge13))/(.6^{12}.4^{13})\rceil=3680673$ 串，共 $20457889$。于是 $H_{.9}=\log 20457889\approx24.28615$ bits，实际定长编号向上取整为 25 位。

\topic{Tut 8 · 二元 Huffman}
\problem{8.1}{哪组不可能是任何分布的二元 Huffman 码？(a) $\{0,10,11\}$；(b) $\{00,01,10,110\}$；(c) $\{01,10\}$。}
\solve (a) 取概率 $(.5,.3,.2)$，先并后两项为 $.5$，再连根；根左得 $0$，右边两叶得 $10,11$。(b) 有三个长 2 和一个长 3，Kraft 和 $3/4+1/8=7/8$；(c) 两个长 2，和 $1/2$。Huffman 每次合并两项，内部节点都有两子树，因此 Kraft 和为 1，(b,c) 不可能。
\problem{8.2}{对任意 $\epsilon>0$，构造分布使最优平均码长 $L>H(X)+1-\epsilon$。}
\solve 取两符号概率 $(1-\delta,\delta)$，其中 $0<\delta<1/2$。两个码字不能用空串，码 $0,1$ 达到最短长度各 1，故 $L=1$。令 $\delta\downarrow0$，两项 $-\delta\log\delta$ 和 $-(1-\delta)\log(1-\delta)$ 均趋 0，故可取 $\delta$ 使 $H<\epsilon$，代入得 $1>H+1-\epsilon$。
\problem{8.3}{(a) 为 $(.3,.3,.2,.1,.1)$ 构造二元 Huffman 码并求平均长。(b) 找新分布使同一个码的平均码长等于新分布的熵。}
\solve (a) 每轮并最小者：$.1+.1=.2$，$.2+.2=.4$，$.3+.3=.6$，$.4+.6=1$。根左放 $.6$ 的两叶，得 $00,01$；右放 $.4$，得 $10$ 和 $110,111$。平均 $L=2(.3+.3+.2)+3(.1+.1)=2.2$。
\parta{b} 让每符号自信息等于对应长度，即 $p'_i=2^{-\ell_i}$。于是 $p'=(1/4,1/4,1/4,1/8,1/8)$，和为 1，$H=L=3(1/4)2+2(1/8)3=2.25$。
\problem{8.4}{概率 $p_1\ge\dots\ge p_{10}$ 的二元 Huffman 码长为 $\ell_1,\ldots,\ell_{10}$。把末项拆成 $\alpha p_{10},(1-\alpha)p_{10}$，$0\le\alpha\le1$。问新的最优码长。}
\solve 两项都不超过 $p_{10}$，可作为最小对先合并；合回 $p_{10}$ 后，剩下的问题与原树相同。取原来第 10 叶并延展一个分支位给两个孩子，得到一组最优长度 $(\ell_1,\ldots,\ell_9,\ell_{10}+1,\ell_{10}+1)$。新增平均位数为 $\alpha p_{10}+(1-\alpha)p_{10}=p_{10}$。若端点时删掉零概率符号，则回到原码。

\topic{Tut 9 · Shannon 与最优码}
\problem{9.1}{概率 $(1/3,1/3,1/4,1/12)$。(a) 构造 Huffman 码。(b) 证明 $(1,2,3,3)$ 和 $(2,2,2,2)$ 都最优。(c) 用例子说明最优码中的某个长度可大于 Shannon 长度。}
\solve (a,b) 先并 $1/12+1/4=1/3$，余三个同权节点。合一原节点与新节点，取 $0,10,110,111$，长度 $(1,2,3,3)$，平均 $1/3+2/3+3(1/4+1/12)=2$。若改合两原节点，四叶都在第 2 层，平均 2。两种打破并列的方式均执行 Huffman，故都最优。(c) 第三符号 $p=1/4$ 的 Shannon 长是 $\lceil\log4\rceil=2$，第一套却用 3 位。
\problem{9.2}{符号概率为权重 $(1,2,3,4,5,6)$ 分别除以 21。(a) 二元 Huffman，(b) 三元 Huffman，(c) 算两者平均码长。}
\solve (a) 二元并 $1+2=3$、$3+3=6$、$4+5=9$、$6+6=12$、$9+12=21$。依原顺序读叶，例 $1110,1111,110,00,01,10$，长 $(4,4,3,2,2,2)$。
\parta{b} 完整三叉树叶数须奇，补零。并 $0+1+2=3$，再 $3+3+4=10$，最后 $5+6+10=21$。去掉虚叶，一组码是 $201,202,21,22,0,1$。
\parta{c} 二元平均为 $(4+8+9+8+10+12)/21=51/21$ bit/符号；三元为 $(3+6+6+8+5+6)/21=34/21$ 三元位/符号。
\problem{9.3}{求使 $\sum_i p_i\ell_i$ 最小的二元即时码长度。(a) $p=(10,9,8,7,7)/41$；(b) $p_i=.9(.1)^{i-1}$，$i=1,2,\ldots$。}
\solve (a) 合并 $7+7=14$、$8+9=17$、$10+14=24$、$17+24=41$，得长 $(2,2,2,3,3)$，平均 $(2\cdot27+3\cdot14)/41=96/41$。
\parta{b} 取 $0,10,110,1110,\ldots$，长 $i$，平均等于几何分布的均值 $1/.9=10/9$。若首项用至少两位，单项贡献 $\ge1.8>10/9$，可排除。首项用一位后，剩余概率 $.1$ 都在另一分支；扣掉共同首位并归一化，分布又是原几何分布。设最优下确界为 $L_*$，所有符号先花一位，尾部至少贡献 $.1L_*$，故 $L_*\ge1+.1L_*$，得 $L_*\ge10/9$，构造达到。
\problem{9.4}{概率 $p_1\ge\dots\ge p_m>0$，$F_i=\sum_{k<i}p_k$，Shannon 码取 $F_i$ 的前 $\ell_i=\lceil\log(1/p_i)\rceil$ 个二进制小数位。(a) 证明前缀性质与 $H(X)\le L<H(X)+1$。(b) 给 $(1/2,1/4,1/8,1/8)$ 构造码。}
\solve (a) $i<j$ 则 $\ell_i\le\ell_j$，且 $F_j-F_i\ge p_i\ge2^{-\ell_i}$。共享 $\ell_i$ 位前缀等于落在同一宽 $2^{-\ell_i}$ 的半开区间，间距不允许，故前缀性质成立。取整给 $-\log p_i\le\ell_i<-\log p_i+1$，乘 $p_i$ 逐项相加即得长度界。(b) $F=(0,1/2,3/4,7/8)$，二进制为 $.0,.10,.110,.111$；长度 $(1,2,3,3)$，取位得 $0,10,110,111$。

\topic{Tut 10 · 信道容量}
\problem{10.1}{输入输出均为 $\{0,\ldots,10\}$，信道 $Y=X+Z\bmod11$，$Z$ 与 $X$ 独立，在 $1,2,3$ 上均匀。(a) 求容量，(b) 给一个最优输入分布。}
\solve 固定 $X=x$，三个输出是 $x+1,x+2,x+3$ 模 11，各占 $1/3$，故 $H(Y|X)=\log3$。$Y$ 最多 11 种，$H(Y)\le\log11$。取 $X$ 均匀，则 $P(Y=y)=\sum_{z=1}^3(1/3)(1/11)=1/11$，输出熵达上界。故 $C=\log11-\log3=\log(11/3)$，均匀输入最优。
\problem{10.2}{二元输入信道，每次以 $\alpha$ 擦除、以 $\epsilon$ 翻转、以 $1-\alpha-\epsilon$ 正确输出。(a) 求容量，(b) 取 $\alpha=0$，(c) 取 $\epsilon=0$。}
\solve (a) 当 $\alpha<1$，令 $E$ 表示擦除。$Y$ 决定 $E$，而擦除概率不随输入变化，故 $I(X;Y)=I(X;E)+I(X;Y|E)$，首项为 0。擦除时 $Y=e$ 恒定；未擦除概率 $1-\alpha$，翻转条件概率 $q=\epsilon/(1-\alpha)$。该分支的 $H(Y|X,E=0)=\B(q)$，输出熵至多 1，均匀输入做到 1。故 $C=(1-\alpha)[1-\B(q)]$。$\alpha=1$ 时恒擦除，容量 0。(b) $C=1-\B(\epsilon)$。(c) $C=1-\alpha$。
\problem{10.3}{$Y=XZ$，$X,Z$ 是独立二元变量，$P(Z=1)=\alpha$。(a) 求容量及最优输入。(b) 若接收方也能观察 $Z$，再求容量。}
\solve (a) 设 $p=P(X=1)$。$P(Y=1)=\alpha p$，且 $X=0$ 时输出恒零、$X=1$ 时输出是 $Z$，所以 $I=\B(\alpha p)-p\B(\alpha)$。$\alpha=0$ 时容量 0，任意输入均可。其余令 $d=\B(\alpha)/\alpha$。令导数 $\alpha\log[(1-\alpha p)/(\alpha p)]-\B(\alpha)=0$。解得 $t=\alpha p^*=1/(1+2^d)$，$p^*=t/\alpha$。导数随 $p$ 严格下降，近 0 时为正；$0<\alpha<1$ 时，在 $p=1$ 导数为 $\alpha\log[(1-\alpha)/\alpha]-\B(\alpha)=\log(1-\alpha)<0$；$\alpha=1$ 时最优为 $p=1/2$。故 $p^*$ 为最大点。由 $d=\log[(1-t)/t]$，代回 $C=\B(t)-td=-\log(1-t)=\log(1+2^{-d})$。
\parta{b} 链式法则 $I(X;Y,Z)=I(X;Z)+I(X;Y|Z)$。独立使第一项为 0；$Z=0$ 不透露输入，$Z=1$ 时 $Y=X$。故为 $\alpha H(X)$，均匀输入给 $C=\alpha$。
\problem{10.4}{证明给信道转移矩阵增加一列不会降低容量。}
\solve 每列是一个输入的输出分布。给新增输入概率 0，其余输入概率按原最优分布设置。所有旧输入和输出的联合概率照旧，故可达到旧容量。新容量是对扩大后的全部输入分布取最大值，因而 $C_{\rm new}\ge C_{\rm old}$。

\topic{Tut 11 · 容量、错误率与译码}
\form{本组共用的码率推导} 均匀独立消息位 $M_1,\ldots,M_K$，译为 $\hat M$，第 $i$ 位错率 $e_i$，平均 $e$。误差位 $M_i\oplus\hat M_i$ 加上 $\hat M_i$ 就能反推 $M_i$，故 $H(M_i|\hat M_i)\le\B(e_i)$。链式法则和条件降熵给 $H(M|\hat M)\le\sum_iH(M_i|\hat M_i)$，Jensen 给 $\sum_i\B(e_i)\le K\B(e)$。于是 $K[1-\B(e)]\le I(M;\hat M)$。数据处理给 $I(M;\hat M)\le I(X^N;Y^N)$；无记忆给 $H(Y^N|X^N)=\sum_i H(Y_i|X_i)$，而 $H(Y^N)\le\sum_iH(Y_i)$，故 $I(X^N;Y^N)\le NC$。两式合并得 $R\le C/[1-\B(e)]$（$e<1/2$）。这是渐近可逼近的上界：均匀二元信源在 Hamming 失真 $e$ 下的率失真函数为 $1-\B(e)$，即每消息位可压到接近这么多比特，再用低于 $C$ 的信道码传送。
\problem{11.1}{BSC 翻转率 $q$，输入概率 $(p,1-p)$。(a) 算互信息，证明 $C=1-\B(q)$。(b) 允许 5\% 错误概率时，求可达码率的最紧上界。}
\solve (a) 输出 0 来自输入 0 未翻转或输入 1 翻转，故 $p_Y(0)=p(1-q)+(1-p)q$。条件输出熵总为 $\B(q)$，所以 $I=\B(q+p(1-2q))-\B(q)$。首项至多 1，$p=1/2$ 时达到，故 $C=1-\B(q)$。(b) 按平均消息位错率解，共用推导给 $R\le[1-\B(q)]/[1-\B(.05)]$，分母约 $.713603$。若题中 5\% 指固定整块错误率，渐近上界由强逆定理给 $C$；若指信道翻转率则代入 $q=.05$，容量约 $.713603$。
\problem{11.2}{BEC 擦除率 $q$，输入概率 $(p,1-p)$。(a) 算互信息并证 $C=1-q$。(b,c) 容许错误概率 25\%、50\% 时分别求最紧码率上界。}
\solve (a) 收到 0 或 1 时输入确定；收到 $e$ 时，由于擦除与输入独立，后验不变。因此 $H(X|Y)=q\B(p)$，互信息 $(1-q)\B(p)$ 在 $p=1/2$ 取 $C=1-q$。(b) 平均消息位错率 $.25$ 给 $R\le(1-q)/[1-\B(.25)]\approx(1-q)/.188722$。(c) 允许 $.5$ 时，什么都不传，对均匀位盲猜的平均错误率就为 $.5$，故此条件不给有限码率上界。若百分比指固定整块错误率，两问的渐近上界都是 $C$；若指信道擦除率，容量分别是 $.75$、$.5$。
\problem{11.3}{输入 $a,b,c$，输出 $a,b,c,d$，转移矩阵四行是 $(3,1,0)/4$、$(1,2,1)/4$、$(0,1,2)/4$、$(0,0,1)/4$。输入概率 $(1/4,1/4,1/2)$。(a) 算互信息。(b) 对每个输出用最大 $P(y|x)$ 的输入译码，算平均错误率。(c,d) 码率分别为 $.01$ 和 $100$ bit/次、位误率 1\% 的声明可能吗？}
\solve (a) 矩阵乘输入：$p_Y(a)=3/16+1/16=1/4$；$p_Y(b)=1/16+2/16+1/8=5/16$；$p_Y(c)=1/16+1/4=5/16$；$p_Y(d)=1/8$。第一列非零概率为 $(3/4,1/4)$，后两列均有 $(1/2,1/4,1/4)$。故 $H(Y|X)=\tfrac14\B(1/4)+\tfrac34(1.5)$；$H(Y)=\tfrac14(2)-2(\tfrac5{16})\log(\tfrac5{16})+\tfrac18(3)$。相减得 $I\approx.59598$。
\parta{b} 按每行最大值，$a,b,c,d$ 分别译成 $a,b,c,c$。每种正确事件用先验乘该条件概率，正确率为 $3/16+1/8+1/4+1/8=11/16$，错率 $5/16$。
\parta{c,d} (a) 的输入已达到 $I>.01$，故 $C>.01$，足够长的码可使整块、从而平均位错误率都低于 1\%。另一方面 $C\le H(X)\le\log3$ 对最大化取值仍成立。共用推导给 $R\le\log3/[1-\B(.01)]\approx1.7243$；100 不可能。
\problem{11.4}{独立加性噪声信道 $Y=X+Z$，$X\in\{0,1\}$，噪声等概率取 $0,a$，$a$ 为任意整数。求容量。}
\solve 给定输入 0、1 时，输出集合分别为 $\{0,a\},\{1,1+a\}$。交集非空仅当 $a=1$（共享 1）或 $a=-1$（共享 0）。共享输出在两输入下都以 $1/2$ 出现，后验保持先验；其它输出确定输入。于是 $H(X|Y)=\tfrac12H(X)$，互信息最大为 $1/2$。其它 $a$ 时集合不交（$a=0$ 按噪声恒 0 理解），所以 $H(X|Y)=0$，均匀输入给 $C=1$。

\topic{Tut 12 · Hamming 与算术编码}
\problem{12.1}{对 $(7,4)$ Hamming 码，(a) 编码 \bb{0111}；(b) 用题目所指维恩图解码 \bb{1100000}；(c) 给该串的校验综合；(d) 最多保证纠正几位？同样 4 个数据位，要用多少位的重复码才能保证相同能力？}
\solve 原题 PDF 缺图，采用常见偶校验约定：数据位置 $3,5,6,7$，三组位置 $(1,3,5,7)$、$(2,3,6,7)$、$(4,5,6,7)$。
\parta{a} 数据位依次 $0,1,1,1$。第一组其余位异或 $0\oplus1\oplus1=0$，第二组 $0\oplus1\oplus1=0$，第三组 $1\oplus1\oplus1=1$，故校验位 $(p_1,p_2,p_4)=(0,0,1)$，组装得 \bb{0001111}。
\parta{b,c} 对 \bb{1100000} 按三组求实际奇偶，得到 $(s_1,s_2,s_4)=(1,1,0)$。综合按高位在前为 \bb{011}，即错误位置 $1+2=3$。翻第 3 位得 \bb{1110000}，取数据位得 \bb{1000}。
\parta{d} 校验矩阵的七列是不同非零的三位位置编号，单列或两列异或都不为零，所以非零码字至少三位；编号 $1,2,3$ 的列异或为零，距离可达 3。半径 1 的译码球互斥，半径 2 的可能重叠，保证纠一位。重复两次遇一错会平票，三次多数决可纠一错，发 $4\cdot3=12$ 位。
\problem{12.2}{首位 0、1 各半。之后下一个符号与前一位相同的概率 $3/5$，不同的概率 $1/5$，结束 EOS 概率 $1/5$。(a) 算术编码 \bb{011}；(b) 解回；(c) 自选一串编码并解回。}
\solve 按 $0,1,\code$ 排序，首轮只有 $0,1$。
\parta{a} 首位 $0$ 取 $[0,.5)$。前位为 0 时，$1$ 的累积段是 $[.6,.8)$，映到 $[.3,.4)$。前位成了 1，$1$ 的累积段是 $[.2,.8)$，映到 $[.32,.38)$。再取 EOS 的最后 20\%，区间 $[.32+.06(.8),.38)=[.368,.38)$。$[95/256,96/256)$ 在其中，所以发 95 的 8 位二进制 \bb{01011111}。
\parta{b} 代表值 $r=95/256=.37109375$。全区间 $r<.5$，读 0；在 $[0,.5)$ 中，$r$ 落在 1 段 $[.3,.4)$；该区间中落在 1 段 $[.32,.38)$；再落在 EOS 段 $[.368,.38)$，读出 \bb{011}。
\parta{c} 选 \bb{0}。首符号取 $[0,.5)$，EOS 切出 $[.4,.5)$。取 $[7/16,8/16)$，发 \bb{0111}。$r=7/16$ 先落 0 段，再落 $[.4,.5)$ 的 EOS 段，译为 \bb{0}。
\problem{12.3}{(a) 为何使用自适应分布？(b) Dirichlet 分布的直观含义？(c) 符号 $a,b,c$ 的虚拟计数为 $(3,1,1)$，EOS 固定 $1/4$。以自适应算术码编码 \bb{bb}。(d) 解回。}
\solve (a) 源的概率未知或会变化时，按观察更新，常见符号获得较宽的区间；编解码同步使用同一更新规则。(b) Dirichlet 是对概率向量的分布。参数可看作虚拟计数；每观察某符号，其计数加 1，预测概率是该计数占总计数的比例。
\parta{c} 选顺序 $a,b,c,\code$。前三项的计数比乘 $3/4$。开始 $a,b,c$ 的概率为 $.45,.15,.15$，故第一个 $b$ 取 $[.45,.60)$。更新 $(3,2,1)$ 后，$a,b$ 概率 $.375,.25$；第二个 $b$ 取 $[.45+.15(.375),.45+.15(.625))=[.50625,.54375)$。更新 $(3,3,1)$；EOS 仍是末 $1/4$，起点 $.50625+.0375(.75)=.534375$，终点 $.54375$。选 $[137/256,138/256)$，发 \bb{10001001}。
\parta{d} $r=137/256=.53515625$。首轮在 $[.45,.60)$，译 $b$ 并更新 $(3,2,1)$；第二轮在 $[.50625,.54375)$，再译 $b$，更新 $(3,3,1)$；第三轮在 $[.534375,.54375)$，读 EOS 停止，输出 \bb{bb}。
\end{multicols}
\end{document}
'
\documentclass[a4paper]{article}
\usepackage[UTF8,scheme=plain,fontset=fandol]{ctex}
\usepackage{amsmath,amssymb,multicol,xcolor}
\ifdefined\LargePrint
  \usepackage[margin=14mm]{geometry}
  \setlength{\columnsep}{6mm}
  \newcommand{\bodyformat}{\fontsize{9.6}{13}\selectfont}
  \newcommand{\cols}{2}
  \setlength{\parskip}{2.7pt}
\else
  \usepackage[top=5.6mm,bottom=7.6mm,left=5.8mm,right=5.8mm]{geometry}
  \setlength{\columnsep}{3mm}
  \newcommand{\bodyformat}{\fontsize{6.65}{7.8}\selectfont}
  \newcommand{\cols}{4}
  \setlength{\parskip}{1.05pt}
\fi
\usepackage{fancyhdr}
\definecolor{ink}{HTML}{163D51}
\pagestyle{fancy}
\fancyhf{}
\fancyfoot[L]{\scriptsize COMP2610 / COMP6261}
\fancyfoot[R]{\scriptsize \thepage}
\renewcommand{\headrulewidth}{0pt}
\setlength{\footskip}{4.5mm}
\setlength{\parindent}{0pt}
\setlength{\multicolsep}{2pt}
\setlength{\emergencystretch}{1.3em}
\setlength{\columnseprule}{.14pt}
\setlength{\premulticols}{0pt}
\setlength{\postmulticols}{0pt}
\raggedcolumns
\newcommand{\topic}[1]{\par\addvspace{2.6pt}{\color{ink}\bfseries #1}\par\nobreak}
\newcommand{\problem}[2]{\par\addvspace{1.4pt}\textcolor{ink}{\textbf{#1}}\enspace\textbf{题}\enspace #2\par\nobreak}
\newcommand{\solve}{\noindent\textbf{解}\enspace}
\newcommand{\parta}[1]{\par\noindent\textbf{(#1)}\enspace}
\newcommand{\form}[1]{\par\noindent\textbf{#1}\enspace}
\newcommand{\E}{\mathbb E}
\newcommand{\V}{\operatorname{Var}}
\newcommand{\B}{h_2}
\newcommand{\bb}[1]{\texttt{#1}}
\newcommand{\code}{\mathrm{EOS}}
\begin{document}
{\color{ink}\fontsize{13}{15}\selectfont\bfseries COMP2610 / COMP6261 \quad L22 公式与 Tut 逐题解}\par
{\fontsize{7}{8.4}\selectfont Tut 1–12 共 46 道去重题。题干以中文保留条件和小问；Tut 5 Q4 与 Tut 7 Q2 共用一份解答。信息量对数以 2 为底。\par}
\begin{multicols}{\cols}
\bodyformat\raggedright
\topic{公式区}
\form{概率与 Bayes} $p(x,y)=p(x)p(y|x)$，$p(y)=\sum_xp(x,y)$，$p(x|y)=p(x,y)/p(y)$。独立时 $p(x,y)=p(x)p(y)$；条件独立在固定条件后用同一乘积形式。
\form{Expected Value} $\E g(X)=\sum_x p(x)g(x)$。期望可直接相加。
\form{Variance} $\V X=\E[X^2]-(\E X)^2$。独立变量相加时方差相加。
\form{Binomial Distribution} $K=\sum_{i=1}^n B_i$，$B_i$ 独立且为 $\mathrm{Bern}(p)$。$P(K=k)=\binom nk p^k(1-p)^{n-k}$，$\E K=np$，$\V K=np(1-p)$。理由：$B_i^2=B_i$，所以单次期望 $p$、方差 $p-p^2$。
\form{Markov's Inequality} $W\ge0,\ a>0$ 时，$P(W\ge a)\le\E W/a$。阈值以上每个结果至少贡献 $a$。
\form{Chebyshev's Inequality} $\mu=\E W,\ t>0$ 时，$P(|W-\mu|\ge t)\le\V(W)/t^2$。单侧阈值 $a>\mu$ 用 $t=a-\mu$。这是对 $(W-\mu)^2$ 用 Markov。
\form{Jensen / MLE} 凸 $f$ 满足 $f(\E X)\le\E f(X)$。i.i.d. 的 MLE 最大化 $\ln L(\theta)=\sum_i\ln p_\theta(x_i)$。
\form{Entropy} 自信息 $-\log p(x)$；$H(X)=-\sum_xp(x)\log p(x)$；$H(X,Y)=-\sum_{x,y}p(x,y)\log p(x,y)$。$H(Y|X)=\sum_xp(x)H(Y|X=x)$。$0\log0=0$。$H(X,Y)=H(X)+H(Y|X)$；$H(X)\le\log|\mathcal X|$。记 $\B(p)=-p\log p-(1-p)\log(1-p)$。
\form{互信息与 KL} $I(X;Y)=H(X)-H(X|Y)=H(Y)-H(Y|X)=D(p_{XY}\Vert p_Xp_Y)$；$D(p\Vert q)=\sum p\log(p/q)\ge0$。条件互信息是给定条件时互信息的加权平均。对 $X\to Y\to Z$，$I(X;Z)\le\min\{I(X;Y),I(Y;Z)\}$。
\form{Extended Ensemble} $X^N$ 把整块视作一个结果。i.i.d. 时 $p(x^N)=\prod_i p(x_i)$，$H(X^N)=NH(X)$。含 $k$ 个 1 的单串概率为 $p^k(1-p)^{N-k}$，同组有 $\binom Nk$ 串。
\form{AEP／典型集} $T_{N,\beta}=\{x^N:|-\log p(x^N)/N-H(X)|<\beta\}$，i.i.d. 时其概率趋 1。最小 $\delta$-充分集按单串概率降序取至累计概率 $\ge\delta$；$H_\delta=\log|S_\delta|$，定长需 $\lceil H_\delta\rceil$ 位。固定 $0<\delta<1$ 时 $H_\delta/N\to H(X)$。
\form{信源编码} $L=\sum_i p_i\ell_i$；二元 Kraft 和 $\sum_i2^{-\ell_i}\le1$。Huffman 每轮并两个最小权重。Shannon 取累计 $F_i=\sum_{j<i}p_j$ 的前 $\ell_i=\lceil-\log p_i\rceil$ 位；SFE 取 $F_i+p_i/2$ 的前 $\ell_i+1$ 位。$H\le L_{\rm opt}\le L_{\rm Shannon}<H+1$。
\form{算术编码} 当前 $[a,a+w)$，符号概率段 $[F,F+p)$，更新为 $[a+wF,a+w(F+p))$。含 EOS 的最终区间内取一个完整二进制网格区间。
\form{信道编码} $Q_{y,x}=p(y|x)$，$p_Y=Qp_X$；$H(Y|X)=\sum_xp(x)H(Q_{\cdot,x})$；$C=\max_{p_X}I(X;Y)$。$R=K/N$；$P_B=P(\hat{\mathbf s}\ne\mathbf s)$；$R<C$ 可使块错率趋零，$R>C$ 不可。

\topic{Tut 1 · 基础概率}
\problem{1.1}{等分转盘标 $1,2,3,4,5$，求 (a) 转到 4，(b) 转到偶数，(c) 转到质数的概率。}
\solve 五格各占 $1/5$，互斥格可加。(a) $\{4\}$：$1/5$。(b) $\{2,4\}$：$2/5$。(c) $\{2,3,5\}$：$3/5$。
\problem{1.2}{ACT 车牌是三字母接三数字。求随机车牌以 ACT 开头且以 7 结尾的概率。}
\solve 按各合法组合等可能且允许重复，总数 $26^3\cdot10^3$。固定前三字母和末位后，中间两位各有 10 种，共 $10^2$。概率 $10^2/(26^3\cdot10^3)=1/175760$。
\problem{1.3}{四条路的测速启用率 $(.4,.3,.2,.3)$，Arya 选路概率 $(.2,.1,.5,.2)$，且总超速。求 (a) 一早被罚，(b) 五早都不被罚。}
\solve (a) 对每条路先乘选路概率和该路处罚率，再相加，得 $.2(.4)+.1(.3)+.5(.2)+.2(.3)=.27$。(b) 每早不被罚概率 $1-.27=.73$。按各早独立且规则不变，五早全不罚为 $.73^5\approx.20731$。
\problem{1.4}{袋内 5 蓝、3 红、2 黄，抽三球。求红球少于 2 颗的概率，分别按 (a) 放回，(b) 不放回。}
\solve (a) 红概率 $.3$，非红 $.7$，三次独立。零红 $.7^3$；一红的出现位置有 3 个，概率 $3(.3)(.7)^2$；和为 $.784$。(b) 把无序三球子集当等可能结果。零红选三非红，有 $\binom73=35$ 种；一红选一红两非红，有 $\binom31\binom72=63$ 种。总数 $\binom{10}3=120$，故 $98/120=49/60$。
\problem{1.5}{两台录像机独立成功，概率分别为 $.7,.6$。求 (a) 零份，(b) 一份，(c) 两份录像的概率。}
\solve 失败率分别是 $.3,.4$。(a) 两失败：$.3(.4)=.12$。(b) 仅第一台或仅第二台，两情形互斥：$.7(.4)+.3(.6)=.46$。(c) 两成功：$.7(.6)=.42$。三项和为 1。

\topic{Tut 2 · 条件概率与 Bayes}
\problem{2.1}{红、蓝、绿箱被选概率为 $(.2,.2,.6)$，各箱苹果、橙子、青柠数分别为 $(3,4,3)$、$(1,1,0)$、$(3,3,4)$。选箱后均匀抽一果。求 (a) 苹果概率，(b) 给定橙子，来自绿箱的概率。}
\solve (a) 各箱苹果比例 $3/10,1/2,3/10$；乘箱概率后总和为 $.06+.10+.18=.34$。(b) 橙的三项联合概率为 $.2(4/10)=.08$、$.2(1/2)=.10$、$.6(3/10)=.18$。用绿箱项除总和，得 $.18/.36=1/2$。
\problem{2.2}{共 1000 次，行 $Y=0,1$、列 $X=0,1$ 的计数为 $(100,250)$、$(150,500)$。求 (a) $p(1,1)$，(b) $P(X=1)$，(c) $\E X$，(d) $P(Y=1|X=1)$，(e) $P(Y=1|X=0)$。(f) 若 $X\ne Y$ 时 $P(Z=1|X,Y)=.9$，相同则为 $.1$，求 $P(X=1,Y=1|Z=1)$。}
\solve (a) 单格 $500/1000=.5$。(b) $X=1$ 列和 $(250+500)/1000=.75$。(c) $0(.25)+1(.75)=.75$。(d) 固定 $X=1$，分母 750，得 $500/750=2/3$。(e) 固定 $X=0$，分母 250，得 $150/250=.6$。(f) 不同格合计 $400/1000=.4$。全概率 $P(Z=1)=.9(.4)+.1(.6)=.42$。目标格且 $Z=1$ 的联合概率是 $.5(.1)=.05$，所以后验 $=.05/.42=5/42$。
\problem{2.3}{独立抛公平硬币两次，每次正面放白球，反面放红球。随后从这两球中放回抽三次，均为红。求箱内两球都红的后验概率。}
\solve 令红球数 $K$。由两次独立抛币，$K=0,1,2$ 的先验为 $(1,2,1)/4$。固定 $K=k$ 后，放回抽每次红率 $k/2$，三红似然为 $(k/2)^3$。乘先验得未归一化后验 $(0,1/16,1/4)$。故 $P(K=2|\text{三红})=(1/4)/(1/16+1/4)=4/5$。
\problem{2.4}{(a) 从条件概率的定义证明 $p(x|y,z)=p(y|x,z)p(x|z)/p(y|z)$。(b) 证明给定 $Y$ 时仍可对任一变量 $X_i$ 求和边缘化。}
\solve (a) 在分母为正时，$p(x|y,z)=p(x,y|z)/p(y|z)$。分子按乘法规则为 $p(x|z)p(y|x,z)$，代回即得。(b) 令 $U$ 代表除 $X_i$ 外的变量。$p(u|y)=p(u,y)/p(y)=[\sum_x p(u,x,y)]/p(y)=\sum_x p(u,x|y)$。

\topic{Tut 3 · MLE 与熵}
\problem{3.1}{取值 $1,2,3$ 的概率为 $(\theta/2,\theta/2,1-\theta)$，$0\le\theta\le1$。i.i.d. 样本中各频数是 $n_1,n_2,n_3$。(a) 写似然。(b) 样本 $3,3,1,2,3,2,2,1,3,1$ 求 MLE。}
\solve (a) 把单次概率相乘，同类因子合并，得 $L=(\theta/2)^{n_1+n_2}(1-\theta)^{n_3}$。(b) 频数 $(3,3,4)$；$\ln L=6\ln(\theta/2)+4\ln(1-\theta)$。令导数 $6/\theta-4/(1-\theta)=0$，得 $6(1-\theta)=4\theta$，即 $\theta=.6$。导数在此点左右分别正负，两端似然为 0，故为最大值。
\problem{3.2}{联合表的行是 $Y=1,2,3,4$，列是 $X=1,2,3,4$，矩阵按 $1/16$ 为单位，四行为 $(0,0,2,2)$、$(2,1,1,0)$、$(2,2,0,0)$、$(0,1,1,2)$。(a) 证不独立。(b) 求 $H(X)$、$H(Y)$、$H(X|Y)$、$H(Y|X)$、$H(X,Y)$。}
\solve (a) 各行各列都合计 $4/16=1/4$。左上为 0，而独立时应为 $1/4\cdot1/4=1/16$。(b) 两个边缘均匀，两个熵都是 $4(1/4)\log 4=2$。每行除以 $1/4$，非零条件概率为 $(1/2,1/2)$ 或 $(1/2,1/4,1/4)$，熵分别为 $1$ 或 $1/2+1/2+1/2=1.5$。四行熵 $(1,1.5,1,1.5)$，均值 $1.25=H(X|Y)$。四列熵 $(1,1.5,1.5,1)$，均值 $1.25=H(Y|X)$。联合表有六个 $1/8$ 和四个 $1/16$，所以 $H(X,Y)=6(1/8)3+4(1/16)4=3.25$。
\problem{3.3}{均匀抽 52 张牌之一。$V,S$ 为点数和花色，人头牌是 J、Q、K。求 (a) 看到红 K 的自信息，(b) 已知是人头牌时看到 K 的条件自信息，(c) $H(S)$ 与 $H(V,S)$。}
\solve (a) 红 K 有两张，$-\log(2/52)=\log26$。(b) 12 张人头牌中四张 K，$-\log(4/12)=\log3$。(c) 花色各有 $13/52=1/4$，$H(S)=2$；一对 $(V,S)$ 确定一张牌，52 对等可能，$H(V,S)=\log52$。
\problem{3.4}{有限离散 $X$。比较 $H(X)$ 与 $H(Y)$，分别取 (a) $Y=2^X$，(b) $Y=\cos X$。}
\solve 确定映射给 $H(Y|X)=0$，两种顺序展开联合熵得 $H(X)=H(Y)+H(X|Y)$。(a) $X=\log_2Y$，故 $H(X|Y)=0$，两熵相等。(b) $H(X|Y)\ge0$，故 $H(Y)\le H(X)$。若 $X$ 等概率为 $\pm1$，余弦相同，熵从 1 变 0；若在支持集上可逆则取等。

\topic{Tut 4 · 熵、Jensen、互信息}
\problem{4.1}{几何变量 $Y$ 的概率为 $p(1-p)^{y-1}$，$y=1,2,\dots$。求期望和方差。}
\solve 设 $q=1-p$、$m=\E Y$、$s=\E Y^2$。首轮成功时 $Y=1$；失败时 $Y=1+Y'$，$Y'$ 同分布。故 $m=p+q(1+m)=1+qm$，得 $m=1/p$。平方后 $s=p+q(1+2m+s)=1+2qm+qs$，故 $s=(1+2q/p)/p=(2-p)/p^2$；方差 $s-m^2=(1-p)/p^2$。
\problem{4.2}{A、B 两队独立且各半胜率，先赢四局者获胜。$X$ 是完整胜负序列，$Y$ 是总局数。求 (a) $H(X)$，(b) $H(Y)$，(c) $H(Y|X)$，(d) $H(X|Y)$。}
\solve 选总冠军有两种；总长 $y$ 时第 $y$ 局必由冠军赢，前 $y-1$ 局选出其前三胜，故 $M_y=2\binom{y-1}3$。$y=4,5,6,7$ 的序列数为 $2,8,20,40$。每序列概率 $2^{-y}$，所以 $p_Y=(1/8,1/4,5/16,5/16)$。
\parta{a} 序列自信息为 $y$，故 $H(X)=4/8+5/4+6(5/16)+7(5/16)=5.8125$。
\parta{b} $H(Y)=(1/8)3+(1/4)2-2(5/16)\log(5/16)\approx1.92379$。
\parta{c,d} 从序列可数出长度，所以 $H(Y|X)=0$。两种链式展开给 $H(X|Y)=H(X)-H(Y)\approx3.88871$。
\problem{4.3}{已知 $\V X=\E[X^2]-(\E X)^2$。用 Jensen 证明方差非负。}
\solve $f(t)=t^2$ 的二阶导数为 $2>0$，故凸。Jensen 给 $f(\E X)\le\E f(X)$，于是 $(\E X)^2\le\E[X^2]$，移项即方差 $\ge0$。严格凸给等号条件为 $X$ 几乎必然恒定。
\problem{4.4}{$X,Y$ 是独立公平 Bernoulli。(a) 算 $I(X;Y)$。(b) 令 $Z=X+Y$，算 $I(X;Y|Z)$。(c) 是否违反数据处理不等式？}
\solve (a) 独立给 $H(X|Y)=H(X)=1$，所以 $I=0$。(b) 四个输入对各占 $1/4$。$Z=0,2$ 时 $X$ 已确定；$Z=1$ 概率 $1/2$，内部两对等可能，故 $H(X|Z)=\tfrac12\cdot1$。从 $X=Z-Y$ 得 $H(X|Y,Z)=0$，故条件互信息 $1/2$。(c) 固定 $Y=0$ 则 $Z=X$，所以给定 $Y$ 后 $X,Z$ 仍相关，$X\to Y\to Z$ 的前提不成立。

\topic{Tut 5 · 基本与概率不等式}
\problem{5.1}{证明 (a) $H(X,Y|Z)\ge H(X|Z)$；(b) $I(X,Y;Z)\ge I(X;Z)$；(c) $H(X,Y,Z)-H(X,Y)\le H(X,Z)-H(X)$，并找取等条件；(d) 证明 $I(X;Z|Y)=I(Z;Y|X)-I(Z;Y)+I(X;Z)$。}
\solve (a) 按条件熵链式展开，左减右为 $H(Y|X,Z)\ge0$；等号要求给定 $(X,Z)$ 后能确定 $Y$。(b) 按互信息链式展开，差为 $I(Y;Z|X)\ge0$；等号要求给定 $X$ 后独立。(c) 左、右分别是 $H(Z|X,Y)$、$H(Z|X)$；右减左就是 $I(Y;Z|X)$，取等同 (b)。(d) $I(X,Y;Z)=I(X;Z)+I(Y;Z|X)=I(Y;Z)+I(X;Z|Y)$；移项再用互信息对称性即得。
\problem{5.2}{比较 (a) $H(5X)$ 与 $H(X)$；(b) $I(g(X);Y)$ 与 $I(X;Y)$；(c) $H(X_0|X_{-1})$ 与 $H(X_0|X_{-1},X_1)$；(d) $H(X,Y)/(H(X)+H(Y))$ 与 $1$。}
\solve (a) 乘、除 5 互逆，故熵相等。(b) $Y\to X\to g(X)$，DPI 给前者 $\le$。(c) 前减后为 $I(X_0;X_1|X_{-1})\ge0$，故前者 $\ge$。(d) 分母减分子是 $I(X;Y)\ge0$，分母为正时比值 $\le1$；两熵都为零时比值无定义。
\problem{5.3}{给定概率质量函数 $p(x)$，证明 $P(p(X)\le d)\log(1/d)\le H(X)$，题面写 $d\ge0$。}
\solve 取 $W=\log(1/p(X))\ge0$，其期望是 $H(X)$。$0<d<1$ 时，$p(X)\le d$ 等价于 $W\ge\log(1/d)>0$。用 Markov 得 $P(p(X)\le d)\le H(X)/\log(1/d)$，乘回阈值。$d\ge1$ 左端非正。$d=0$ 原式无定义；有限熵时，写 $t=\log(1/d)$，有 $tP(W\ge t)\le\E[W1_{\{W\ge t\}}]\to0$，可按极限补为 0。
\form{5.4} 与 7.2 相同，共用下方完整题解。

\topic{Tut 6 · 熵分组、KL、Markov 链}
\problem{6.1}{分布 $p=(p_1,\ldots,p_m)$，把末两项合并成 $q$，新末项为 $s=p_{m-1}+p_m$。证明 $H(p)=H(q)+sH(p_{m-1}/s,p_m/s)$。}
\solve 设 $Y$ 为原结果 $X$ 的组标签。$H(Y|X)=0$，链式得 $H(X)=H(Y)+H(X|Y)$。$H(Y)=H(q)$。未合并的组内只有一个结果，条件熵为 0；合并组用原概率除以组和 $s$，该组熵乘出现概率 $s$，得到右边。若 $s=0$，该组贡献按 0 计。
\problem{6.2}{对联合分布 $p(x,y,z)$，证明 (a) $I(X,Y;Z)-I(Y,Z;X)=I(Y;Z)-I(X;Y)$；(b) $D(p_{XYZ}\Vert p_Xp_Yp_Z)=I(X;Y)+I(X,Y;Z)$。}
\solve (a) 两项依次展开为 $I(Y;Z)+I(X;Z|Y)$ 与 $I(Y;X)+I(Z;X|Y)$。条件项相等，相减抵消。
\parta{b} 在联合概率正的格子，$p(x,y,z)/[p(x)p(y)p(z)]$ 等于 $p(x,y,z)/[p(x,y)p(z)]$ 乘 $p(x,y)/[p(x)p(y)]$。取对数后按联合分布求和，第一项就是 $I(X,Y;Z)$。第二项对数与 $z$ 无关，先算 $\sum_zp(x,y,z)=p(x,y)$，余项就是 $I(X;Y)$。
\problem{6.3}{马尔可夫链 $X_1\to X_2\to X_3$ 的状态数依次为 $n,k,m$，其中 $k<n,k<m$。(a) 求 $I(X_1;X_3)$ 的上界。(b) $k=1$ 时求其值及独立性。}
\solve (a) 数据处理给 $I(X_1;X_3)\le I(X_1;X_2)$。后者等于 $H(X_2)-H(X_2|X_1)\le H(X_2)$；与均匀分布的 KL 给 $\log k-H(X_2)\ge0$，故上界 $\log k$。(b) $k=1$ 时上界为 0；互信息非负，所以等于 0，两端独立。
\problem{6.4}{正面概率 $1/5$ 的硬币独立抛偶数 $N$ 次。用 (a) Markov，(b) Chebyshev 给「至少一半正面」求上界。}
\solve 令 $K\sim\mathrm{Bin}(N,1/5)$，$\E K=N/5$，$\V K=N(1/5)(4/5)=4N/25$。(a) 阈值 $N/2$，故 $P(K\ge N/2)\le(N/5)/(N/2)=2/5$。(b) $K\ge N/2$ 蕴含 $|K-N/5|\ge N/2-N/5=3N/10$，所求概率不超过 $(4N/25)/(3N/10)^2=16/(9N)$；概率界可与 1 取较小值。

\topic{Tut 7 · AEP、计数与定长编码}
\problem{7.1}{正面率 $p$ 的硬币独立抛 $n$ 次。(a) 算正面不少于 $n-1$ 次的精确概率，(b) 用 Markov 求上界，(c) $n=2$ 时上界在精确值的 1\% 以内，求 $p$。}
\solve $K\sim\mathrm{Bin}(n,p)$。(a) 只有 $K=n-1,n$，相加得 $np^{n-1}(1-p)+p^n$。(b) $n\ge2$ 时阈值为正，界为 $\E K/(n-1)=np/(n-1)$，也可与 1 取最小；$n=1$ 时事件必然。
\parta{c} $n=2$ 时精确值 $2p-p^2$，原 Markov 界 $2p$，差 $p^2$。按相对误差，$p>0$ 时 $p^2/(2p-p^2)=p/(2-p)\le.01$，得 $101p\le2$；$p=0$ 两值相同。若指 1 个百分点，$p^2\le.01$，得 $p\le.1$。若先截断上界，在 $p\ge1/2$ 段差为 $(1-p)^2$；相对误差 $(1-p)^2/[1-(1-p)^2]\le.01$ 给 $p\ge1-1/\sqrt{101}$；绝对误差 $(1-p)^2\le.01$ 给 $p\ge.9$。
\problem{7.2 = 5.4}{i.i.d. 二元信源 $p(1)=.005$、$p(0)=.995$，每 100 位为一块。仅给至多 3 个 1 的块编码。(a) 定长码至少几位？(b) 未编码块的概率？(c) 用 Chebyshev 给上界并比较。}
\solve 令 $K$ 为块内 1 的数目。
\parta{a} 恰 $k$ 个 1 只需选位置，有 $\binom{100}k$ 串。编码数为 $1+100+4950+161700=166751$。因为 $2^{17}<166751\le2^{18}$，故最短 18 位。
\parta{b} 每个含 $k$ 个 1 的串概率 $.005^k .995^{100-k}$，乘串数给 $P(K=k)$。未编码是 $K\ge4$，概率 $1-\sum_{k=0}^3\binom{100}k(.005)^k(.995)^{100-k}\approx.00167327$。
\parta{c} $\mu=100(.005)=.5$，$\sigma^2=100(.005)(.995)=.4975$。若 $K\ge4$ 则 $|K-.5|\ge3.5$，故 $P(K\ge4)\le.4975/3.5^2\approx.0406122$；上界大于精确值。
% Keep the full typical-set exercise on the second side of the compact edition.
\ifdefined\LargePrint\else
  \end{multicols}
  \newpage
  \begin{multicols}{\cols}
  \bodyformat\raggedright
\fi
\problem{7.3}{二元信源 $(p(0),p(1))=(.4,.6)$。(a) 求 $H(X)$。(b) 扩展长度 $N=25$，$\beta=.1$。求 (1) 哪些串典型，(2) 典型集合总概率，(3) 典型串数，(4) 最小 $.9$-充分集合大小，(5) 其 essential bit content。}
\solve (a) $H=-.4\log .4-.6\log .6\approx.97095$。
\parta{b1} 若有 $k$ 个 1，单串概率 $.6^k .4^{25-k}$。每符号自信息为 $-(k/25)\log .6-(1-k/25)\log .4$；减去 $H$ 并合并对数，得 $(.6-k/25)\log(1.5)$。要求绝对值小于 $.1$，即 $15-2.5/\log(1.5)<k<15+2.5/\log(1.5)$；边界约 $10.73,19.27$，所以 $k=11,\ldots,19$。
\parta{b2,b3} 每组有 $\binom{25}k$ 串。乘单串概率再加组得 $\sum_{k=11}^{19}\binom{25}k(.6)^k(.4)^{25-k}\approx.936246$；只加各组串数得 $26366510$。
\parta{b4,b5} $k$ 每加 1，单串概率乘 $1.5$；贪心取最高概率串，任何低概率入选串可换成未选的高概率串来增大覆盖率。先取 $k\ge13$ 的 $\sum_{k=13}^{25}\binom{25}k=16777216$ 串，覆盖率 $\sum_{k=13}^{25}\binom{25}k(.6)^k(.4)^{25-k}\approx.84623223$。再取 $k=12$ 的 $\lceil(.9-P(K\ge13))/(.6^{12}.4^{13})\rceil=3680673$ 串，共 $20457889$。于是 $H_{.9}=\log 20457889\approx24.28615$ bits，实际定长编号向上取整为 25 位。

\topic{Tut 8 · 二元 Huffman}
\problem{8.1}{哪组不可能是任何分布的二元 Huffman 码？(a) $\{0,10,11\}$；(b) $\{00,01,10,110\}$；(c) $\{01,10\}$。}
\solve (a) 取概率 $(.5,.3,.2)$，先并后两项为 $.5$，再连根；根左得 $0$，右边两叶得 $10,11$。(b) 有三个长 2 和一个长 3，Kraft 和 $3/4+1/8=7/8$；(c) 两个长 2，和 $1/2$。Huffman 每次合并两项，内部节点都有两子树，因此 Kraft 和为 1，(b,c) 不可能。
\problem{8.2}{对任意 $\epsilon>0$，构造分布使最优平均码长 $L>H(X)+1-\epsilon$。}
\solve 取两符号概率 $(1-\delta,\delta)$，其中 $0<\delta<1/2$。两个码字不能用空串，码 $0,1$ 达到最短长度各 1，故 $L=1$。令 $\delta\downarrow0$，两项 $-\delta\log\delta$ 和 $-(1-\delta)\log(1-\delta)$ 均趋 0，故可取 $\delta$ 使 $H<\epsilon$，代入得 $1>H+1-\epsilon$。
\problem{8.3}{(a) 为 $(.3,.3,.2,.1,.1)$ 构造二元 Huffman 码并求平均长。(b) 找新分布使同一个码的平均码长等于新分布的熵。}
\solve (a) 每轮并最小者：$.1+.1=.2$，$.2+.2=.4$，$.3+.3=.6$，$.4+.6=1$。根左放 $.6$ 的两叶，得 $00,01$；右放 $.4$，得 $10$ 和 $110,111$。平均 $L=2(.3+.3+.2)+3(.1+.1)=2.2$。
\parta{b} 让每符号自信息等于对应长度，即 $p'_i=2^{-\ell_i}$。于是 $p'=(1/4,1/4,1/4,1/8,1/8)$，和为 1，$H=L=3(1/4)2+2(1/8)3=2.25$。
\problem{8.4}{概率 $p_1\ge\dots\ge p_{10}$ 的二元 Huffman 码长为 $\ell_1,\ldots,\ell_{10}$。把末项拆成 $\alpha p_{10},(1-\alpha)p_{10}$，$0\le\alpha\le1$。问新的最优码长。}
\solve 两项都不超过 $p_{10}$，可作为最小对先合并；合回 $p_{10}$ 后，剩下的问题与原树相同。取原来第 10 叶并延展一个分支位给两个孩子，得到一组最优长度 $(\ell_1,\ldots,\ell_9,\ell_{10}+1,\ell_{10}+1)$。新增平均位数为 $\alpha p_{10}+(1-\alpha)p_{10}=p_{10}$。若端点时删掉零概率符号，则回到原码。

\topic{Tut 9 · Shannon 与最优码}
\problem{9.1}{概率 $(1/3,1/3,1/4,1/12)$。(a) 构造 Huffman 码。(b) 证明 $(1,2,3,3)$ 和 $(2,2,2,2)$ 都最优。(c) 用例子说明最优码中的某个长度可大于 Shannon 长度。}
\solve (a,b) 先并 $1/12+1/4=1/3$，余三个同权节点。合一原节点与新节点，取 $0,10,110,111$，长度 $(1,2,3,3)$，平均 $1/3+2/3+3(1/4+1/12)=2$。若改合两原节点，四叶都在第 2 层，平均 2。两种打破并列的方式均执行 Huffman，故都最优。(c) 第三符号 $p=1/4$ 的 Shannon 长是 $\lceil\log4\rceil=2$，第一套却用 3 位。
\problem{9.2}{符号概率为权重 $(1,2,3,4,5,6)$ 分别除以 21。(a) 二元 Huffman，(b) 三元 Huffman，(c) 算两者平均码长。}
\solve (a) 二元并 $1+2=3$、$3+3=6$、$4+5=9$、$6+6=12$、$9+12=21$。依原顺序读叶，例 $1110,1111,110,00,01,10$，长 $(4,4,3,2,2,2)$。
\parta{b} 完整三叉树叶数须奇，补零。并 $0+1+2=3$，再 $3+3+4=10$，最后 $5+6+10=21$。去掉虚叶，一组码是 $201,202,21,22,0,1$。
\parta{c} 二元平均为 $(4+8+9+8+10+12)/21=51/21$ bit/符号；三元为 $(3+6+6+8+5+6)/21=34/21$ 三元位/符号。
\problem{9.3}{求使 $\sum_i p_i\ell_i$ 最小的二元即时码长度。(a) $p=(10,9,8,7,7)/41$；(b) $p_i=.9(.1)^{i-1}$，$i=1,2,\ldots$。}
\solve (a) 合并 $7+7=14$、$8+9=17$、$10+14=24$、$17+24=41$，得长 $(2,2,2,3,3)$，平均 $(2\cdot27+3\cdot14)/41=96/41$。
\parta{b} 取 $0,10,110,1110,\ldots$，长 $i$，平均等于几何分布的均值 $1/.9=10/9$。若首项用至少两位，单项贡献 $\ge1.8>10/9$，可排除。首项用一位后，剩余概率 $.1$ 都在另一分支；扣掉共同首位并归一化，分布又是原几何分布。设最优下确界为 $L_*$，所有符号先花一位，尾部至少贡献 $.1L_*$，故 $L_*\ge1+.1L_*$，得 $L_*\ge10/9$，构造达到。
\problem{9.4}{概率 $p_1\ge\dots\ge p_m>0$，$F_i=\sum_{k<i}p_k$，Shannon 码取 $F_i$ 的前 $\ell_i=\lceil\log(1/p_i)\rceil$ 个二进制小数位。(a) 证明前缀性质与 $H(X)\le L<H(X)+1$。(b) 给 $(1/2,1/4,1/8,1/8)$ 构造码。}
\solve (a) $i<j$ 则 $\ell_i\le\ell_j$，且 $F_j-F_i\ge p_i\ge2^{-\ell_i}$。共享 $\ell_i$ 位前缀等于落在同一宽 $2^{-\ell_i}$ 的半开区间，间距不允许，故前缀性质成立。取整给 $-\log p_i\le\ell_i<-\log p_i+1$，乘 $p_i$ 逐项相加即得长度界。(b) $F=(0,1/2,3/4,7/8)$，二进制为 $.0,.10,.110,.111$；长度 $(1,2,3,3)$，取位得 $0,10,110,111$。

\topic{Tut 10 · 信道容量}
\problem{10.1}{输入输出均为 $\{0,\ldots,10\}$，信道 $Y=X+Z\bmod11$，$Z$ 与 $X$ 独立，在 $1,2,3$ 上均匀。(a) 求容量，(b) 给一个最优输入分布。}
\solve 固定 $X=x$，三个输出是 $x+1,x+2,x+3$ 模 11，各占 $1/3$，故 $H(Y|X)=\log3$。$Y$ 最多 11 种，$H(Y)\le\log11$。取 $X$ 均匀，则 $P(Y=y)=\sum_{z=1}^3(1/3)(1/11)=1/11$，输出熵达上界。故 $C=\log11-\log3=\log(11/3)$，均匀输入最优。
\problem{10.2}{二元输入信道，每次以 $\alpha$ 擦除、以 $\epsilon$ 翻转、以 $1-\alpha-\epsilon$ 正确输出。(a) 求容量，(b) 取 $\alpha=0$，(c) 取 $\epsilon=0$。}
\solve (a) 当 $\alpha<1$，令 $E$ 表示擦除。$Y$ 决定 $E$，而擦除概率不随输入变化，故 $I(X;Y)=I(X;E)+I(X;Y|E)$，首项为 0。擦除时 $Y=e$ 恒定；未擦除概率 $1-\alpha$，翻转条件概率 $q=\epsilon/(1-\alpha)$。该分支的 $H(Y|X,E=0)=\B(q)$，输出熵至多 1，均匀输入做到 1。故 $C=(1-\alpha)[1-\B(q)]$。$\alpha=1$ 时恒擦除，容量 0。(b) $C=1-\B(\epsilon)$。(c) $C=1-\alpha$。
\problem{10.3}{$Y=XZ$，$X,Z$ 是独立二元变量，$P(Z=1)=\alpha$。(a) 求容量及最优输入。(b) 若接收方也能观察 $Z$，再求容量。}
\solve (a) 设 $p=P(X=1)$。$P(Y=1)=\alpha p$，且 $X=0$ 时输出恒零、$X=1$ 时输出是 $Z$，所以 $I=\B(\alpha p)-p\B(\alpha)$。$\alpha=0$ 时容量 0，任意输入均可。其余令 $d=\B(\alpha)/\alpha$。令导数 $\alpha\log[(1-\alpha p)/(\alpha p)]-\B(\alpha)=0$。解得 $t=\alpha p^*=1/(1+2^d)$，$p^*=t/\alpha$。导数随 $p$ 严格下降，近 0 时为正；$0<\alpha<1$ 时，在 $p=1$ 导数为 $\alpha\log[(1-\alpha)/\alpha]-\B(\alpha)=\log(1-\alpha)<0$；$\alpha=1$ 时最优为 $p=1/2$。故 $p^*$ 为最大点。由 $d=\log[(1-t)/t]$，代回 $C=\B(t)-td=-\log(1-t)=\log(1+2^{-d})$。
\parta{b} 链式法则 $I(X;Y,Z)=I(X;Z)+I(X;Y|Z)$。独立使第一项为 0；$Z=0$ 不透露输入，$Z=1$ 时 $Y=X$。故为 $\alpha H(X)$，均匀输入给 $C=\alpha$。
\problem{10.4}{证明给信道转移矩阵增加一列不会降低容量。}
\solve 每列是一个输入的输出分布。给新增输入概率 0，其余输入概率按原最优分布设置。所有旧输入和输出的联合概率照旧，故可达到旧容量。新容量是对扩大后的全部输入分布取最大值，因而 $C_{\rm new}\ge C_{\rm old}$。

\topic{Tut 11 · 容量、错误率与译码}
\form{本组共用的码率推导} 均匀独立消息位 $M_1,\ldots,M_K$，译为 $\hat M$，第 $i$ 位错率 $e_i$，平均 $e$。误差位 $M_i\oplus\hat M_i$ 加上 $\hat M_i$ 就能反推 $M_i$，故 $H(M_i|\hat M_i)\le\B(e_i)$。链式法则和条件降熵给 $H(M|\hat M)\le\sum_iH(M_i|\hat M_i)$，Jensen 给 $\sum_i\B(e_i)\le K\B(e)$。于是 $K[1-\B(e)]\le I(M;\hat M)$。数据处理给 $I(M;\hat M)\le I(X^N;Y^N)$；无记忆给 $H(Y^N|X^N)=\sum_i H(Y_i|X_i)$，而 $H(Y^N)\le\sum_iH(Y_i)$，故 $I(X^N;Y^N)\le NC$。两式合并得 $R\le C/[1-\B(e)]$（$e<1/2$）。这是渐近可逼近的上界：均匀二元信源在 Hamming 失真 $e$ 下的率失真函数为 $1-\B(e)$，即每消息位可压到接近这么多比特，再用低于 $C$ 的信道码传送。
\problem{11.1}{BSC 翻转率 $q$，输入概率 $(p,1-p)$。(a) 算互信息，证明 $C=1-\B(q)$。(b) 允许 5\% 错误概率时，求可达码率的最紧上界。}
\solve (a) 输出 0 来自输入 0 未翻转或输入 1 翻转，故 $p_Y(0)=p(1-q)+(1-p)q$。条件输出熵总为 $\B(q)$，所以 $I=\B(q+p(1-2q))-\B(q)$。首项至多 1，$p=1/2$ 时达到，故 $C=1-\B(q)$。(b) 按平均消息位错率解，共用推导给 $R\le[1-\B(q)]/[1-\B(.05)]$，分母约 $.713603$。若题中 5\% 指固定整块错误率，渐近上界由强逆定理给 $C$；若指信道翻转率则代入 $q=.05$，容量约 $.713603$。
\problem{11.2}{BEC 擦除率 $q$，输入概率 $(p,1-p)$。(a) 算互信息并证 $C=1-q$。(b,c) 容许错误概率 25\%、50\% 时分别求最紧码率上界。}
\solve (a) 收到 0 或 1 时输入确定；收到 $e$ 时，由于擦除与输入独立，后验不变。因此 $H(X|Y)=q\B(p)$，互信息 $(1-q)\B(p)$ 在 $p=1/2$ 取 $C=1-q$。(b) 平均消息位错率 $.25$ 给 $R\le(1-q)/[1-\B(.25)]\approx(1-q)/.188722$。(c) 允许 $.5$ 时，什么都不传，对均匀位盲猜的平均错误率就为 $.5$，故此条件不给有限码率上界。若百分比指固定整块错误率，两问的渐近上界都是 $C$；若指信道擦除率，容量分别是 $.75$、$.5$。
\problem{11.3}{输入 $a,b,c$，输出 $a,b,c,d$，转移矩阵四行是 $(3,1,0)/4$、$(1,2,1)/4$、$(0,1,2)/4$、$(0,0,1)/4$。输入概率 $(1/4,1/4,1/2)$。(a) 算互信息。(b) 对每个输出用最大 $P(y|x)$ 的输入译码，算平均错误率。(c,d) 码率分别为 $.01$ 和 $100$ bit/次、位误率 1\% 的声明可能吗？}
\solve (a) 矩阵乘输入：$p_Y(a)=3/16+1/16=1/4$；$p_Y(b)=1/16+2/16+1/8=5/16$；$p_Y(c)=1/16+1/4=5/16$；$p_Y(d)=1/8$。第一列非零概率为 $(3/4,1/4)$，后两列均有 $(1/2,1/4,1/4)$。故 $H(Y|X)=\tfrac14\B(1/4)+\tfrac34(1.5)$；$H(Y)=\tfrac14(2)-2(\tfrac5{16})\log(\tfrac5{16})+\tfrac18(3)$。相减得 $I\approx.59598$。
\parta{b} 按每行最大值，$a,b,c,d$ 分别译成 $a,b,c,c$。每种正确事件用先验乘该条件概率，正确率为 $3/16+1/8+1/4+1/8=11/16$，错率 $5/16$。
\parta{c,d} (a) 的输入已达到 $I>.01$，故 $C>.01$，足够长的码可使整块、从而平均位错误率都低于 1\%。另一方面 $C\le H(X)\le\log3$ 对最大化取值仍成立。共用推导给 $R\le\log3/[1-\B(.01)]\approx1.7243$；100 不可能。
\problem{11.4}{独立加性噪声信道 $Y=X+Z$，$X\in\{0,1\}$，噪声等概率取 $0,a$，$a$ 为任意整数。求容量。}
\solve 给定输入 0、1 时，输出集合分别为 $\{0,a\},\{1,1+a\}$。交集非空仅当 $a=1$（共享 1）或 $a=-1$（共享 0）。共享输出在两输入下都以 $1/2$ 出现，后验保持先验；其它输出确定输入。于是 $H(X|Y)=\tfrac12H(X)$，互信息最大为 $1/2$。其它 $a$ 时集合不交（$a=0$ 按噪声恒 0 理解），所以 $H(X|Y)=0$，均匀输入给 $C=1$。

\topic{Tut 12 · Hamming 与算术编码}
\problem{12.1}{对 $(7,4)$ Hamming 码，(a) 编码 \bb{0111}；(b) 用题目所指维恩图解码 \bb{1100000}；(c) 给该串的校验综合；(d) 最多保证纠正几位？同样 4 个数据位，要用多少位的重复码才能保证相同能力？}
\solve 原题 PDF 缺图，采用常见偶校验约定：数据位置 $3,5,6,7$，三组位置 $(1,3,5,7)$、$(2,3,6,7)$、$(4,5,6,7)$。
\parta{a} 数据位依次 $0,1,1,1$。第一组其余位异或 $0\oplus1\oplus1=0$，第二组 $0\oplus1\oplus1=0$，第三组 $1\oplus1\oplus1=1$，故校验位 $(p_1,p_2,p_4)=(0,0,1)$，组装得 \bb{0001111}。
\parta{b,c} 对 \bb{1100000} 按三组求实际奇偶，得到 $(s_1,s_2,s_4)=(1,1,0)$。综合按高位在前为 \bb{011}，即错误位置 $1+2=3$。翻第 3 位得 \bb{1110000}，取数据位得 \bb{1000}。
\parta{d} 校验矩阵的七列是不同非零的三位位置编号，单列或两列异或都不为零，所以非零码字至少三位；编号 $1,2,3$ 的列异或为零，距离可达 3。半径 1 的译码球互斥，半径 2 的可能重叠，保证纠一位。重复两次遇一错会平票，三次多数决可纠一错，发 $4\cdot3=12$ 位。
\problem{12.2}{首位 0、1 各半。之后下一个符号与前一位相同的概率 $3/5$，不同的概率 $1/5$，结束 EOS 概率 $1/5$。(a) 算术编码 \bb{011}；(b) 解回；(c) 自选一串编码并解回。}
\solve 按 $0,1,\code$ 排序，首轮只有 $0,1$。
\parta{a} 首位 $0$ 取 $[0,.5)$。前位为 0 时，$1$ 的累积段是 $[.6,.8)$，映到 $[.3,.4)$。前位成了 1，$1$ 的累积段是 $[.2,.8)$，映到 $[.32,.38)$。再取 EOS 的最后 20\%，区间 $[.32+.06(.8),.38)=[.368,.38)$。$[95/256,96/256)$ 在其中，所以发 95 的 8 位二进制 \bb{01011111}。
\parta{b} 代表值 $r=95/256=.37109375$。全区间 $r<.5$，读 0；在 $[0,.5)$ 中，$r$ 落在 1 段 $[.3,.4)$；该区间中落在 1 段 $[.32,.38)$；再落在 EOS 段 $[.368,.38)$，读出 \bb{011}。
\parta{c} 选 \bb{0}。首符号取 $[0,.5)$，EOS 切出 $[.4,.5)$。取 $[7/16,8/16)$，发 \bb{0111}。$r=7/16$ 先落 0 段，再落 $[.4,.5)$ 的 EOS 段，译为 \bb{0}。
\problem{12.3}{(a) 为何使用自适应分布？(b) Dirichlet 分布的直观含义？(c) 符号 $a,b,c$ 的虚拟计数为 $(3,1,1)$，EOS 固定 $1/4$。以自适应算术码编码 \bb{bb}。(d) 解回。}
\solve (a) 源的概率未知或会变化时，按观察更新，常见符号获得较宽的区间；编解码同步使用同一更新规则。(b) Dirichlet 是对概率向量的分布。参数可看作虚拟计数；每观察某符号，其计数加 1，预测概率是该计数占总计数的比例。
\parta{c} 选顺序 $a,b,c,\code$。前三项的计数比乘 $3/4$。开始 $a,b,c$ 的概率为 $.45,.15,.15$，故第一个 $b$ 取 $[.45,.60)$。更新 $(3,2,1)$ 后，$a,b$ 概率 $.375,.25$；第二个 $b$ 取 $[.45+.15(.375),.45+.15(.625))=[.50625,.54375)$。更新 $(3,3,1)$；EOS 仍是末 $1/4$，起点 $.50625+.0375(.75)=.534375$，终点 $.54375$。选 $[137/256,138/256)$，发 \bb{10001001}。
\parta{d} $r=137/256=.53515625$。首轮在 $[.45,.60)$，译 $b$ 并更新 $(3,2,1)$；第二轮在 $[.50625,.54375)$，再译 $b$，更新 $(3,3,1)$；第三轮在 $[.534375,.54375)$，读 EOS 停止，输出 \bb{bb}。
\end{multicols}
\end{document}
'
\documentclass[a4paper]{article}
\usepackage[UTF8,scheme=plain,fontset=fandol]{ctex}
\usepackage{amsmath,amssymb,multicol,xcolor}
\ifdefined\LargePrint
  \usepackage[margin=14mm]{geometry}
  \setlength{\columnsep}{6mm}
  \newcommand{\bodyformat}{\fontsize{9.6}{13}\selectfont}
  \newcommand{\cols}{2}
  \setlength{\parskip}{2.7pt}
\else
  \usepackage[top=5.6mm,bottom=7.6mm,left=5.8mm,right=5.8mm]{geometry}
  \setlength{\columnsep}{3mm}
  \newcommand{\bodyformat}{\fontsize{6.65}{7.8}\selectfont}
  \newcommand{\cols}{4}
  \setlength{\parskip}{1.05pt}
\fi
\usepackage{fancyhdr}
\definecolor{ink}{HTML}{163D51}
\pagestyle{fancy}
\fancyhf{}
\fancyfoot[L]{\scriptsize COMP2610 / COMP6261}
\fancyfoot[R]{\scriptsize \thepage}
\renewcommand{\headrulewidth}{0pt}
\setlength{\footskip}{4.5mm}
\setlength{\parindent}{0pt}
\setlength{\multicolsep}{2pt}
\setlength{\emergencystretch}{1.3em}
\setlength{\columnseprule}{.14pt}
\setlength{\premulticols}{0pt}
\setlength{\postmulticols}{0pt}
\raggedcolumns
\newcommand{\topic}[1]{\par\addvspace{2.6pt}{\color{ink}\bfseries #1}\par\nobreak}
\newcommand{\problem}[2]{\par\addvspace{1.4pt}\textcolor{ink}{\textbf{#1}}\enspace\textbf{题}\enspace #2\par\nobreak}
\newcommand{\solve}{\noindent\textbf{解}\enspace}
\newcommand{\parta}[1]{\par\noindent\textbf{(#1)}\enspace}
\newcommand{\form}[1]{\par\noindent\textbf{#1}\enspace}
\newcommand{\E}{\mathbb E}
\newcommand{\V}{\operatorname{Var}}
\newcommand{\B}{h_2}
\newcommand{\bb}[1]{\texttt{#1}}
\newcommand{\code}{\mathrm{EOS}}
\begin{document}
{\color{ink}\fontsize{13}{15}\selectfont\bfseries COMP2610 / COMP6261 \quad L22 公式与 Tut 逐题解}\par
{\fontsize{7}{8.4}\selectfont Tut 1–12 共 46 道去重题。题干以中文保留条件和小问；Tut 5 Q4 与 Tut 7 Q2 共用一份解答。信息量对数以 2 为底。\par}
\begin{multicols}{\cols}
\bodyformat\raggedright
\topic{公式区}
\form{概率与 Bayes} $p(x,y)=p(x)p(y|x)$，$p(y)=\sum_xp(x,y)$，$p(x|y)=p(x,y)/p(y)$。独立时 $p(x,y)=p(x)p(y)$；条件独立在固定条件后用同一乘积形式。
\form{Expected Value} $\E g(X)=\sum_x p(x)g(x)$。期望可直接相加。
\form{Variance} $\V X=\E[X^2]-(\E X)^2$。独立变量相加时方差相加。
\form{Binomial Distribution} $K=\sum_{i=1}^n B_i$，$B_i$ 独立且为 $\mathrm{Bern}(p)$。$P(K=k)=\binom nk p^k(1-p)^{n-k}$，$\E K=np$，$\V K=np(1-p)$。理由：$B_i^2=B_i$，所以单次期望 $p$、方差 $p-p^2$。
\form{Markov's Inequality} $W\ge0,\ a>0$ 时，$P(W\ge a)\le\E W/a$。阈值以上每个结果至少贡献 $a$。
\form{Chebyshev's Inequality} $\mu=\E W,\ t>0$ 时，$P(|W-\mu|\ge t)\le\V(W)/t^2$。单侧阈值 $a>\mu$ 用 $t=a-\mu$。这是对 $(W-\mu)^2$ 用 Markov。
\form{Jensen / MLE} 凸 $f$ 满足 $f(\E X)\le\E f(X)$。i.i.d. 的 MLE 最大化 $\ln L(\theta)=\sum_i\ln p_\theta(x_i)$。
\form{Entropy} 自信息 $-\log p(x)$；$H(X)=-\sum_xp(x)\log p(x)$；$H(X,Y)=-\sum_{x,y}p(x,y)\log p(x,y)$。$H(Y|X)=\sum_xp(x)H(Y|X=x)$。$0\log0=0$。$H(X,Y)=H(X)+H(Y|X)$；$H(X)\le\log|\mathcal X|$。记 $\B(p)=-p\log p-(1-p)\log(1-p)$。
\form{互信息与 KL} $I(X;Y)=H(X)-H(X|Y)=H(Y)-H(Y|X)=D(p_{XY}\Vert p_Xp_Y)$；$D(p\Vert q)=\sum p\log(p/q)\ge0$。条件互信息是给定条件时互信息的加权平均。对 $X\to Y\to Z$，$I(X;Z)\le\min\{I(X;Y),I(Y;Z)\}$。
\form{Extended Ensemble} $X^N$ 把整块视作一个结果。i.i.d. 时 $p(x^N)=\prod_i p(x_i)$，$H(X^N)=NH(X)$。含 $k$ 个 1 的单串概率为 $p^k(1-p)^{N-k}$，同组有 $\binom Nk$ 串。
\form{AEP／典型集} $T_{N,\beta}=\{x^N:|-\log p(x^N)/N-H(X)|<\beta\}$，i.i.d. 时其概率趋 1。最小 $\delta$-充分集按单串概率降序取至累计概率 $\ge\delta$；$H_\delta=\log|S_\delta|$，定长需 $\lceil H_\delta\rceil$ 位。固定 $0<\delta<1$ 时 $H_\delta/N\to H(X)$。
\form{信源编码} $L=\sum_i p_i\ell_i$；二元 Kraft 和 $\sum_i2^{-\ell_i}\le1$。Huffman 每轮并两个最小权重。Shannon 取累计 $F_i=\sum_{j<i}p_j$ 的前 $\ell_i=\lceil-\log p_i\rceil$ 位；SFE 取 $F_i+p_i/2$ 的前 $\ell_i+1$ 位。$H\le L_{\rm opt}\le L_{\rm Shannon}<H+1$。
\form{算术编码} 当前 $[a,a+w)$，符号概率段 $[F,F+p)$，更新为 $[a+wF,a+w(F+p))$。含 EOS 的最终区间内取一个完整二进制网格区间。
\form{信道编码} $Q_{y,x}=p(y|x)$，$p_Y=Qp_X$；$H(Y|X)=\sum_xp(x)H(Q_{\cdot,x})$；$C=\max_{p_X}I(X;Y)$。$R=K/N$；$P_B=P(\hat{\mathbf s}\ne\mathbf s)$；$R<C$ 可使块错率趋零，$R>C$ 不可。

\topic{Tut 1 · 基础概率}
\problem{1.1}{等分转盘标 $1,2,3,4,5$，求 (a) 转到 4，(b) 转到偶数，(c) 转到质数的概率。}
\solve 五格各占 $1/5$，互斥格可加。(a) $\{4\}$：$1/5$。(b) $\{2,4\}$：$2/5$。(c) $\{2,3,5\}$：$3/5$。
\problem{1.2}{ACT 车牌是三字母接三数字。求随机车牌以 ACT 开头且以 7 结尾的概率。}
\solve 按各合法组合等可能且允许重复，总数 $26^3\cdot10^3$。固定前三字母和末位后，中间两位各有 10 种，共 $10^2$。概率 $10^2/(26^3\cdot10^3)=1/175760$。
\problem{1.3}{四条路的测速启用率 $(.4,.3,.2,.3)$，Arya 选路概率 $(.2,.1,.5,.2)$，且总超速。求 (a) 一早被罚，(b) 五早都不被罚。}
\solve (a) 对每条路先乘选路概率和该路处罚率，再相加，得 $.2(.4)+.1(.3)+.5(.2)+.2(.3)=.27$。(b) 每早不被罚概率 $1-.27=.73$。按各早独立且规则不变，五早全不罚为 $.73^5\approx.20731$。
\problem{1.4}{袋内 5 蓝、3 红、2 黄，抽三球。求红球少于 2 颗的概率，分别按 (a) 放回，(b) 不放回。}
\solve (a) 红概率 $.3$，非红 $.7$，三次独立。零红 $.7^3$；一红的出现位置有 3 个，概率 $3(.3)(.7)^2$；和为 $.784$。(b) 把无序三球子集当等可能结果。零红选三非红，有 $\binom73=35$ 种；一红选一红两非红，有 $\binom31\binom72=63$ 种。总数 $\binom{10}3=120$，故 $98/120=49/60$。
\problem{1.5}{两台录像机独立成功，概率分别为 $.7,.6$。求 (a) 零份，(b) 一份，(c) 两份录像的概率。}
\solve 失败率分别是 $.3,.4$。(a) 两失败：$.3(.4)=.12$。(b) 仅第一台或仅第二台，两情形互斥：$.7(.4)+.3(.6)=.46$。(c) 两成功：$.7(.6)=.42$。三项和为 1。

\topic{Tut 2 · 条件概率与 Bayes}
\problem{2.1}{红、蓝、绿箱被选概率为 $(.2,.2,.6)$，各箱苹果、橙子、青柠数分别为 $(3,4,3)$、$(1,1,0)$、$(3,3,4)$。选箱后均匀抽一果。求 (a) 苹果概率，(b) 给定橙子，来自绿箱的概率。}
\solve (a) 各箱苹果比例 $3/10,1/2,3/10$；乘箱概率后总和为 $.06+.10+.18=.34$。(b) 橙的三项联合概率为 $.2(4/10)=.08$、$.2(1/2)=.10$、$.6(3/10)=.18$。用绿箱项除总和，得 $.18/.36=1/2$。
\problem{2.2}{共 1000 次，行 $Y=0,1$、列 $X=0,1$ 的计数为 $(100,250)$、$(150,500)$。求 (a) $p(1,1)$，(b) $P(X=1)$，(c) $\E X$，(d) $P(Y=1|X=1)$，(e) $P(Y=1|X=0)$。(f) 若 $X\ne Y$ 时 $P(Z=1|X,Y)=.9$，相同则为 $.1$，求 $P(X=1,Y=1|Z=1)$。}
\solve (a) 单格 $500/1000=.5$。(b) $X=1$ 列和 $(250+500)/1000=.75$。(c) $0(.25)+1(.75)=.75$。(d) 固定 $X=1$，分母 750，得 $500/750=2/3$。(e) 固定 $X=0$，分母 250，得 $150/250=.6$。(f) 不同格合计 $400/1000=.4$。全概率 $P(Z=1)=.9(.4)+.1(.6)=.42$。目标格且 $Z=1$ 的联合概率是 $.5(.1)=.05$，所以后验 $=.05/.42=5/42$。
\problem{2.3}{独立抛公平硬币两次，每次正面放白球，反面放红球。随后从这两球中放回抽三次，均为红。求箱内两球都红的后验概率。}
\solve 令红球数 $K$。由两次独立抛币，$K=0,1,2$ 的先验为 $(1,2,1)/4$。固定 $K=k$ 后，放回抽每次红率 $k/2$，三红似然为 $(k/2)^3$。乘先验得未归一化后验 $(0,1/16,1/4)$。故 $P(K=2|\text{三红})=(1/4)/(1/16+1/4)=4/5$。
\problem{2.4}{(a) 从条件概率的定义证明 $p(x|y,z)=p(y|x,z)p(x|z)/p(y|z)$。(b) 证明给定 $Y$ 时仍可对任一变量 $X_i$ 求和边缘化。}
\solve (a) 在分母为正时，$p(x|y,z)=p(x,y|z)/p(y|z)$。分子按乘法规则为 $p(x|z)p(y|x,z)$，代回即得。(b) 令 $U$ 代表除 $X_i$ 外的变量。$p(u|y)=p(u,y)/p(y)=[\sum_x p(u,x,y)]/p(y)=\sum_x p(u,x|y)$。

\topic{Tut 3 · MLE 与熵}
\problem{3.1}{取值 $1,2,3$ 的概率为 $(\theta/2,\theta/2,1-\theta)$，$0\le\theta\le1$。i.i.d. 样本中各频数是 $n_1,n_2,n_3$。(a) 写似然。(b) 样本 $3,3,1,2,3,2,2,1,3,1$ 求 MLE。}
\solve (a) 把单次概率相乘，同类因子合并，得 $L=(\theta/2)^{n_1+n_2}(1-\theta)^{n_3}$。(b) 频数 $(3,3,4)$；$\ln L=6\ln(\theta/2)+4\ln(1-\theta)$。令导数 $6/\theta-4/(1-\theta)=0$，得 $6(1-\theta)=4\theta$，即 $\theta=.6$。导数在此点左右分别正负，两端似然为 0，故为最大值。
\problem{3.2}{联合表的行是 $Y=1,2,3,4$，列是 $X=1,2,3,4$，矩阵按 $1/16$ 为单位，四行为 $(0,0,2,2)$、$(2,1,1,0)$、$(2,2,0,0)$、$(0,1,1,2)$。(a) 证不独立。(b) 求 $H(X)$、$H(Y)$、$H(X|Y)$、$H(Y|X)$、$H(X,Y)$。}
\solve (a) 各行各列都合计 $4/16=1/4$。左上为 0，而独立时应为 $1/4\cdot1/4=1/16$。(b) 两个边缘均匀，两个熵都是 $4(1/4)\log 4=2$。每行除以 $1/4$，非零条件概率为 $(1/2,1/2)$ 或 $(1/2,1/4,1/4)$，熵分别为 $1$ 或 $1/2+1/2+1/2=1.5$。四行熵 $(1,1.5,1,1.5)$，均值 $1.25=H(X|Y)$。四列熵 $(1,1.5,1.5,1)$，均值 $1.25=H(Y|X)$。联合表有六个 $1/8$ 和四个 $1/16$，所以 $H(X,Y)=6(1/8)3+4(1/16)4=3.25$。
\problem{3.3}{均匀抽 52 张牌之一。$V,S$ 为点数和花色，人头牌是 J、Q、K。求 (a) 看到红 K 的自信息，(b) 已知是人头牌时看到 K 的条件自信息，(c) $H(S)$ 与 $H(V,S)$。}
\solve (a) 红 K 有两张，$-\log(2/52)=\log26$。(b) 12 张人头牌中四张 K，$-\log(4/12)=\log3$。(c) 花色各有 $13/52=1/4$，$H(S)=2$；一对 $(V,S)$ 确定一张牌，52 对等可能，$H(V,S)=\log52$。
\problem{3.4}{有限离散 $X$。比较 $H(X)$ 与 $H(Y)$，分别取 (a) $Y=2^X$，(b) $Y=\cos X$。}
\solve 确定映射给 $H(Y|X)=0$，两种顺序展开联合熵得 $H(X)=H(Y)+H(X|Y)$。(a) $X=\log_2Y$，故 $H(X|Y)=0$，两熵相等。(b) $H(X|Y)\ge0$，故 $H(Y)\le H(X)$。若 $X$ 等概率为 $\pm1$，余弦相同，熵从 1 变 0；若在支持集上可逆则取等。

\topic{Tut 4 · 熵、Jensen、互信息}
\problem{4.1}{几何变量 $Y$ 的概率为 $p(1-p)^{y-1}$，$y=1,2,\dots$。求期望和方差。}
\solve 设 $q=1-p$、$m=\E Y$、$s=\E Y^2$。首轮成功时 $Y=1$；失败时 $Y=1+Y'$，$Y'$ 同分布。故 $m=p+q(1+m)=1+qm$，得 $m=1/p$。平方后 $s=p+q(1+2m+s)=1+2qm+qs$，故 $s=(1+2q/p)/p=(2-p)/p^2$；方差 $s-m^2=(1-p)/p^2$。
\problem{4.2}{A、B 两队独立且各半胜率，先赢四局者获胜。$X$ 是完整胜负序列，$Y$ 是总局数。求 (a) $H(X)$，(b) $H(Y)$，(c) $H(Y|X)$，(d) $H(X|Y)$。}
\solve 选总冠军有两种；总长 $y$ 时第 $y$ 局必由冠军赢，前 $y-1$ 局选出其前三胜，故 $M_y=2\binom{y-1}3$。$y=4,5,6,7$ 的序列数为 $2,8,20,40$。每序列概率 $2^{-y}$，所以 $p_Y=(1/8,1/4,5/16,5/16)$。
\parta{a} 序列自信息为 $y$，故 $H(X)=4/8+5/4+6(5/16)+7(5/16)=5.8125$。
\parta{b} $H(Y)=(1/8)3+(1/4)2-2(5/16)\log(5/16)\approx1.92379$。
\parta{c,d} 从序列可数出长度，所以 $H(Y|X)=0$。两种链式展开给 $H(X|Y)=H(X)-H(Y)\approx3.88871$。
\problem{4.3}{已知 $\V X=\E[X^2]-(\E X)^2$。用 Jensen 证明方差非负。}
\solve $f(t)=t^2$ 的二阶导数为 $2>0$，故凸。Jensen 给 $f(\E X)\le\E f(X)$，于是 $(\E X)^2\le\E[X^2]$，移项即方差 $\ge0$。严格凸给等号条件为 $X$ 几乎必然恒定。
\problem{4.4}{$X,Y$ 是独立公平 Bernoulli。(a) 算 $I(X;Y)$。(b) 令 $Z=X+Y$，算 $I(X;Y|Z)$。(c) 是否违反数据处理不等式？}
\solve (a) 独立给 $H(X|Y)=H(X)=1$，所以 $I=0$。(b) 四个输入对各占 $1/4$。$Z=0,2$ 时 $X$ 已确定；$Z=1$ 概率 $1/2$，内部两对等可能，故 $H(X|Z)=\tfrac12\cdot1$。从 $X=Z-Y$ 得 $H(X|Y,Z)=0$，故条件互信息 $1/2$。(c) 固定 $Y=0$ 则 $Z=X$，所以给定 $Y$ 后 $X,Z$ 仍相关，$X\to Y\to Z$ 的前提不成立。

\topic{Tut 5 · 基本与概率不等式}
\problem{5.1}{证明 (a) $H(X,Y|Z)\ge H(X|Z)$；(b) $I(X,Y;Z)\ge I(X;Z)$；(c) $H(X,Y,Z)-H(X,Y)\le H(X,Z)-H(X)$，并找取等条件；(d) 证明 $I(X;Z|Y)=I(Z;Y|X)-I(Z;Y)+I(X;Z)$。}
\solve (a) 按条件熵链式展开，左减右为 $H(Y|X,Z)\ge0$；等号要求给定 $(X,Z)$ 后能确定 $Y$。(b) 按互信息链式展开，差为 $I(Y;Z|X)\ge0$；等号要求给定 $X$ 后独立。(c) 左、右分别是 $H(Z|X,Y)$、$H(Z|X)$；右减左就是 $I(Y;Z|X)$，取等同 (b)。(d) $I(X,Y;Z)=I(X;Z)+I(Y;Z|X)=I(Y;Z)+I(X;Z|Y)$；移项再用互信息对称性即得。
\problem{5.2}{比较 (a) $H(5X)$ 与 $H(X)$；(b) $I(g(X);Y)$ 与 $I(X;Y)$；(c) $H(X_0|X_{-1})$ 与 $H(X_0|X_{-1},X_1)$；(d) $H(X,Y)/(H(X)+H(Y))$ 与 $1$。}
\solve (a) 乘、除 5 互逆，故熵相等。(b) $Y\to X\to g(X)$，DPI 给前者 $\le$。(c) 前减后为 $I(X_0;X_1|X_{-1})\ge0$，故前者 $\ge$。(d) 分母减分子是 $I(X;Y)\ge0$，分母为正时比值 $\le1$；两熵都为零时比值无定义。
\problem{5.3}{给定概率质量函数 $p(x)$，证明 $P(p(X)\le d)\log(1/d)\le H(X)$，题面写 $d\ge0$。}
\solve 取 $W=\log(1/p(X))\ge0$，其期望是 $H(X)$。$0<d<1$ 时，$p(X)\le d$ 等价于 $W\ge\log(1/d)>0$。用 Markov 得 $P(p(X)\le d)\le H(X)/\log(1/d)$，乘回阈值。$d\ge1$ 左端非正。$d=0$ 原式无定义；有限熵时，写 $t=\log(1/d)$，有 $tP(W\ge t)\le\E[W1_{\{W\ge t\}}]\to0$，可按极限补为 0。
\form{5.4} 与 7.2 相同，共用下方完整题解。

\topic{Tut 6 · 熵分组、KL、Markov 链}
\problem{6.1}{分布 $p=(p_1,\ldots,p_m)$，把末两项合并成 $q$，新末项为 $s=p_{m-1}+p_m$。证明 $H(p)=H(q)+sH(p_{m-1}/s,p_m/s)$。}
\solve 设 $Y$ 为原结果 $X$ 的组标签。$H(Y|X)=0$，链式得 $H(X)=H(Y)+H(X|Y)$。$H(Y)=H(q)$。未合并的组内只有一个结果，条件熵为 0；合并组用原概率除以组和 $s$，该组熵乘出现概率 $s$，得到右边。若 $s=0$，该组贡献按 0 计。
\problem{6.2}{对联合分布 $p(x,y,z)$，证明 (a) $I(X,Y;Z)-I(Y,Z;X)=I(Y;Z)-I(X;Y)$；(b) $D(p_{XYZ}\Vert p_Xp_Yp_Z)=I(X;Y)+I(X,Y;Z)$。}
\solve (a) 两项依次展开为 $I(Y;Z)+I(X;Z|Y)$ 与 $I(Y;X)+I(Z;X|Y)$。条件项相等，相减抵消。
\parta{b} 在联合概率正的格子，$p(x,y,z)/[p(x)p(y)p(z)]$ 等于 $p(x,y,z)/[p(x,y)p(z)]$ 乘 $p(x,y)/[p(x)p(y)]$。取对数后按联合分布求和，第一项就是 $I(X,Y;Z)$。第二项对数与 $z$ 无关，先算 $\sum_zp(x,y,z)=p(x,y)$，余项就是 $I(X;Y)$。
\problem{6.3}{马尔可夫链 $X_1\to X_2\to X_3$ 的状态数依次为 $n,k,m$，其中 $k<n,k<m$。(a) 求 $I(X_1;X_3)$ 的上界。(b) $k=1$ 时求其值及独立性。}
\solve (a) 数据处理给 $I(X_1;X_3)\le I(X_1;X_2)$。后者等于 $H(X_2)-H(X_2|X_1)\le H(X_2)$；与均匀分布的 KL 给 $\log k-H(X_2)\ge0$，故上界 $\log k$。(b) $k=1$ 时上界为 0；互信息非负，所以等于 0，两端独立。
\problem{6.4}{正面概率 $1/5$ 的硬币独立抛偶数 $N$ 次。用 (a) Markov，(b) Chebyshev 给「至少一半正面」求上界。}
\solve 令 $K\sim\mathrm{Bin}(N,1/5)$，$\E K=N/5$，$\V K=N(1/5)(4/5)=4N/25$。(a) 阈值 $N/2$，故 $P(K\ge N/2)\le(N/5)/(N/2)=2/5$。(b) $K\ge N/2$ 蕴含 $|K-N/5|\ge N/2-N/5=3N/10$，所求概率不超过 $(4N/25)/(3N/10)^2=16/(9N)$；概率界可与 1 取较小值。

\topic{Tut 7 · AEP、计数与定长编码}
\problem{7.1}{正面率 $p$ 的硬币独立抛 $n$ 次。(a) 算正面不少于 $n-1$ 次的精确概率，(b) 用 Markov 求上界，(c) $n=2$ 时上界在精确值的 1\% 以内，求 $p$。}
\solve $K\sim\mathrm{Bin}(n,p)$。(a) 只有 $K=n-1,n$，相加得 $np^{n-1}(1-p)+p^n$。(b) $n\ge2$ 时阈值为正，界为 $\E K/(n-1)=np/(n-1)$，也可与 1 取最小；$n=1$ 时事件必然。
\parta{c} $n=2$ 时精确值 $2p-p^2$，原 Markov 界 $2p$，差 $p^2$。按相对误差，$p>0$ 时 $p^2/(2p-p^2)=p/(2-p)\le.01$，得 $101p\le2$；$p=0$ 两值相同。若指 1 个百分点，$p^2\le.01$，得 $p\le.1$。若先截断上界，在 $p\ge1/2$ 段差为 $(1-p)^2$；相对误差 $(1-p)^2/[1-(1-p)^2]\le.01$ 给 $p\ge1-1/\sqrt{101}$；绝对误差 $(1-p)^2\le.01$ 给 $p\ge.9$。
\problem{7.2 = 5.4}{i.i.d. 二元信源 $p(1)=.005$、$p(0)=.995$，每 100 位为一块。仅给至多 3 个 1 的块编码。(a) 定长码至少几位？(b) 未编码块的概率？(c) 用 Chebyshev 给上界并比较。}
\solve 令 $K$ 为块内 1 的数目。
\parta{a} 恰 $k$ 个 1 只需选位置，有 $\binom{100}k$ 串。编码数为 $1+100+4950+161700=166751$。因为 $2^{17}<166751\le2^{18}$，故最短 18 位。
\parta{b} 每个含 $k$ 个 1 的串概率 $.005^k .995^{100-k}$，乘串数给 $P(K=k)$。未编码是 $K\ge4$，概率 $1-\sum_{k=0}^3\binom{100}k(.005)^k(.995)^{100-k}\approx.00167327$。
\parta{c} $\mu=100(.005)=.5$，$\sigma^2=100(.005)(.995)=.4975$。若 $K\ge4$ 则 $|K-.5|\ge3.5$，故 $P(K\ge4)\le.4975/3.5^2\approx.0406122$；上界大于精确值。
% Keep the full typical-set exercise on the second side of the compact edition.
\ifdefined\LargePrint\else
  \end{multicols}
  \newpage
  \begin{multicols}{\cols}
  \bodyformat\raggedright
\fi
\problem{7.3}{二元信源 $(p(0),p(1))=(.4,.6)$。(a) 求 $H(X)$。(b) 扩展长度 $N=25$，$\beta=.1$。求 (1) 哪些串典型，(2) 典型集合总概率，(3) 典型串数，(4) 最小 $.9$-充分集合大小，(5) 其 essential bit content。}
\solve (a) $H=-.4\log .4-.6\log .6\approx.97095$。
\parta{b1} 若有 $k$ 个 1，单串概率 $.6^k .4^{25-k}$。每符号自信息为 $-(k/25)\log .6-(1-k/25)\log .4$；减去 $H$ 并合并对数，得 $(.6-k/25)\log(1.5)$。要求绝对值小于 $.1$，即 $15-2.5/\log(1.5)<k<15+2.5/\log(1.5)$；边界约 $10.73,19.27$，所以 $k=11,\ldots,19$。
\parta{b2,b3} 每组有 $\binom{25}k$ 串。乘单串概率再加组得 $\sum_{k=11}^{19}\binom{25}k(.6)^k(.4)^{25-k}\approx.936246$；只加各组串数得 $26366510$。
\parta{b4,b5} $k$ 每加 1，单串概率乘 $1.5$；贪心取最高概率串，任何低概率入选串可换成未选的高概率串来增大覆盖率。先取 $k\ge13$ 的 $\sum_{k=13}^{25}\binom{25}k=16777216$ 串，覆盖率 $\sum_{k=13}^{25}\binom{25}k(.6)^k(.4)^{25-k}\approx.84623223$。再取 $k=12$ 的 $\lceil(.9-P(K\ge13))/(.6^{12}.4^{13})\rceil=3680673$ 串，共 $20457889$。于是 $H_{.9}=\log 20457889\approx24.28615$ bits，实际定长编号向上取整为 25 位。

\topic{Tut 8 · 二元 Huffman}
\problem{8.1}{哪组不可能是任何分布的二元 Huffman 码？(a) $\{0,10,11\}$；(b) $\{00,01,10,110\}$；(c) $\{01,10\}$。}
\solve (a) 取概率 $(.5,.3,.2)$，先并后两项为 $.5$，再连根；根左得 $0$，右边两叶得 $10,11$。(b) 有三个长 2 和一个长 3，Kraft 和 $3/4+1/8=7/8$；(c) 两个长 2，和 $1/2$。Huffman 每次合并两项，内部节点都有两子树，因此 Kraft 和为 1，(b,c) 不可能。
\problem{8.2}{对任意 $\epsilon>0$，构造分布使最优平均码长 $L>H(X)+1-\epsilon$。}
\solve 取两符号概率 $(1-\delta,\delta)$，其中 $0<\delta<1/2$。两个码字不能用空串，码 $0,1$ 达到最短长度各 1，故 $L=1$。令 $\delta\downarrow0$，两项 $-\delta\log\delta$ 和 $-(1-\delta)\log(1-\delta)$ 均趋 0，故可取 $\delta$ 使 $H<\epsilon$，代入得 $1>H+1-\epsilon$。
\problem{8.3}{(a) 为 $(.3,.3,.2,.1,.1)$ 构造二元 Huffman 码并求平均长。(b) 找新分布使同一个码的平均码长等于新分布的熵。}
\solve (a) 每轮并最小者：$.1+.1=.2$，$.2+.2=.4$，$.3+.3=.6$，$.4+.6=1$。根左放 $.6$ 的两叶，得 $00,01$；右放 $.4$，得 $10$ 和 $110,111$。平均 $L=2(.3+.3+.2)+3(.1+.1)=2.2$。
\parta{b} 让每符号自信息等于对应长度，即 $p'_i=2^{-\ell_i}$。于是 $p'=(1/4,1/4,1/4,1/8,1/8)$，和为 1，$H=L=3(1/4)2+2(1/8)3=2.25$。
\problem{8.4}{概率 $p_1\ge\dots\ge p_{10}$ 的二元 Huffman 码长为 $\ell_1,\ldots,\ell_{10}$。把末项拆成 $\alpha p_{10},(1-\alpha)p_{10}$，$0\le\alpha\le1$。问新的最优码长。}
\solve 两项都不超过 $p_{10}$，可作为最小对先合并；合回 $p_{10}$ 后，剩下的问题与原树相同。取原来第 10 叶并延展一个分支位给两个孩子，得到一组最优长度 $(\ell_1,\ldots,\ell_9,\ell_{10}+1,\ell_{10}+1)$。新增平均位数为 $\alpha p_{10}+(1-\alpha)p_{10}=p_{10}$。若端点时删掉零概率符号，则回到原码。

\topic{Tut 9 · Shannon 与最优码}
\problem{9.1}{概率 $(1/3,1/3,1/4,1/12)$。(a) 构造 Huffman 码。(b) 证明 $(1,2,3,3)$ 和 $(2,2,2,2)$ 都最优。(c) 用例子说明最优码中的某个长度可大于 Shannon 长度。}
\solve (a,b) 先并 $1/12+1/4=1/3$，余三个同权节点。合一原节点与新节点，取 $0,10,110,111$，长度 $(1,2,3,3)$，平均 $1/3+2/3+3(1/4+1/12)=2$。若改合两原节点，四叶都在第 2 层，平均 2。两种打破并列的方式均执行 Huffman，故都最优。(c) 第三符号 $p=1/4$ 的 Shannon 长是 $\lceil\log4\rceil=2$，第一套却用 3 位。
\problem{9.2}{符号概率为权重 $(1,2,3,4,5,6)$ 分别除以 21。(a) 二元 Huffman，(b) 三元 Huffman，(c) 算两者平均码长。}
\solve (a) 二元并 $1+2=3$、$3+3=6$、$4+5=9$、$6+6=12$、$9+12=21$。依原顺序读叶，例 $1110,1111,110,00,01,10$，长 $(4,4,3,2,2,2)$。
\parta{b} 完整三叉树叶数须奇，补零。并 $0+1+2=3$，再 $3+3+4=10$，最后 $5+6+10=21$。去掉虚叶，一组码是 $201,202,21,22,0,1$。
\parta{c} 二元平均为 $(4+8+9+8+10+12)/21=51/21$ bit/符号；三元为 $(3+6+6+8+5+6)/21=34/21$ 三元位/符号。
\problem{9.3}{求使 $\sum_i p_i\ell_i$ 最小的二元即时码长度。(a) $p=(10,9,8,7,7)/41$；(b) $p_i=.9(.1)^{i-1}$，$i=1,2,\ldots$。}
\solve (a) 合并 $7+7=14$、$8+9=17$、$10+14=24$、$17+24=41$，得长 $(2,2,2,3,3)$，平均 $(2\cdot27+3\cdot14)/41=96/41$。
\parta{b} 取 $0,10,110,1110,\ldots$，长 $i$，平均等于几何分布的均值 $1/.9=10/9$。若首项用至少两位，单项贡献 $\ge1.8>10/9$，可排除。首项用一位后，剩余概率 $.1$ 都在另一分支；扣掉共同首位并归一化，分布又是原几何分布。设最优下确界为 $L_*$，所有符号先花一位，尾部至少贡献 $.1L_*$，故 $L_*\ge1+.1L_*$，得 $L_*\ge10/9$，构造达到。
\problem{9.4}{概率 $p_1\ge\dots\ge p_m>0$，$F_i=\sum_{k<i}p_k$，Shannon 码取 $F_i$ 的前 $\ell_i=\lceil\log(1/p_i)\rceil$ 个二进制小数位。(a) 证明前缀性质与 $H(X)\le L<H(X)+1$。(b) 给 $(1/2,1/4,1/8,1/8)$ 构造码。}
\solve (a) $i<j$ 则 $\ell_i\le\ell_j$，且 $F_j-F_i\ge p_i\ge2^{-\ell_i}$。共享 $\ell_i$ 位前缀等于落在同一宽 $2^{-\ell_i}$ 的半开区间，间距不允许，故前缀性质成立。取整给 $-\log p_i\le\ell_i<-\log p_i+1$，乘 $p_i$ 逐项相加即得长度界。(b) $F=(0,1/2,3/4,7/8)$，二进制为 $.0,.10,.110,.111$；长度 $(1,2,3,3)$，取位得 $0,10,110,111$。

\topic{Tut 10 · 信道容量}
\problem{10.1}{输入输出均为 $\{0,\ldots,10\}$，信道 $Y=X+Z\bmod11$，$Z$ 与 $X$ 独立，在 $1,2,3$ 上均匀。(a) 求容量，(b) 给一个最优输入分布。}
\solve 固定 $X=x$，三个输出是 $x+1,x+2,x+3$ 模 11，各占 $1/3$，故 $H(Y|X)=\log3$。$Y$ 最多 11 种，$H(Y)\le\log11$。取 $X$ 均匀，则 $P(Y=y)=\sum_{z=1}^3(1/3)(1/11)=1/11$，输出熵达上界。故 $C=\log11-\log3=\log(11/3)$，均匀输入最优。
\problem{10.2}{二元输入信道，每次以 $\alpha$ 擦除、以 $\epsilon$ 翻转、以 $1-\alpha-\epsilon$ 正确输出。(a) 求容量，(b) 取 $\alpha=0$，(c) 取 $\epsilon=0$。}
\solve (a) 当 $\alpha<1$，令 $E$ 表示擦除。$Y$ 决定 $E$，而擦除概率不随输入变化，故 $I(X;Y)=I(X;E)+I(X;Y|E)$，首项为 0。擦除时 $Y=e$ 恒定；未擦除概率 $1-\alpha$，翻转条件概率 $q=\epsilon/(1-\alpha)$。该分支的 $H(Y|X,E=0)=\B(q)$，输出熵至多 1，均匀输入做到 1。故 $C=(1-\alpha)[1-\B(q)]$。$\alpha=1$ 时恒擦除，容量 0。(b) $C=1-\B(\epsilon)$。(c) $C=1-\alpha$。
\problem{10.3}{$Y=XZ$，$X,Z$ 是独立二元变量，$P(Z=1)=\alpha$。(a) 求容量及最优输入。(b) 若接收方也能观察 $Z$，再求容量。}
\solve (a) 设 $p=P(X=1)$。$P(Y=1)=\alpha p$，且 $X=0$ 时输出恒零、$X=1$ 时输出是 $Z$，所以 $I=\B(\alpha p)-p\B(\alpha)$。$\alpha=0$ 时容量 0，任意输入均可。其余令 $d=\B(\alpha)/\alpha$。令导数 $\alpha\log[(1-\alpha p)/(\alpha p)]-\B(\alpha)=0$。解得 $t=\alpha p^*=1/(1+2^d)$，$p^*=t/\alpha$。导数随 $p$ 严格下降，近 0 时为正；$0<\alpha<1$ 时，在 $p=1$ 导数为 $\alpha\log[(1-\alpha)/\alpha]-\B(\alpha)=\log(1-\alpha)<0$；$\alpha=1$ 时最优为 $p=1/2$。故 $p^*$ 为最大点。由 $d=\log[(1-t)/t]$，代回 $C=\B(t)-td=-\log(1-t)=\log(1+2^{-d})$。
\parta{b} 链式法则 $I(X;Y,Z)=I(X;Z)+I(X;Y|Z)$。独立使第一项为 0；$Z=0$ 不透露输入，$Z=1$ 时 $Y=X$。故为 $\alpha H(X)$，均匀输入给 $C=\alpha$。
\problem{10.4}{证明给信道转移矩阵增加一列不会降低容量。}
\solve 每列是一个输入的输出分布。给新增输入概率 0，其余输入概率按原最优分布设置。所有旧输入和输出的联合概率照旧，故可达到旧容量。新容量是对扩大后的全部输入分布取最大值，因而 $C_{\rm new}\ge C_{\rm old}$。

\topic{Tut 11 · 容量、错误率与译码}
\form{本组共用的码率推导} 均匀独立消息位 $M_1,\ldots,M_K$，译为 $\hat M$，第 $i$ 位错率 $e_i$，平均 $e$。误差位 $M_i\oplus\hat M_i$ 加上 $\hat M_i$ 就能反推 $M_i$，故 $H(M_i|\hat M_i)\le\B(e_i)$。链式法则和条件降熵给 $H(M|\hat M)\le\sum_iH(M_i|\hat M_i)$，Jensen 给 $\sum_i\B(e_i)\le K\B(e)$。于是 $K[1-\B(e)]\le I(M;\hat M)$。数据处理给 $I(M;\hat M)\le I(X^N;Y^N)$；无记忆给 $H(Y^N|X^N)=\sum_i H(Y_i|X_i)$，而 $H(Y^N)\le\sum_iH(Y_i)$，故 $I(X^N;Y^N)\le NC$。两式合并得 $R\le C/[1-\B(e)]$（$e<1/2$）。这是渐近可逼近的上界：均匀二元信源在 Hamming 失真 $e$ 下的率失真函数为 $1-\B(e)$，即每消息位可压到接近这么多比特，再用低于 $C$ 的信道码传送。
\problem{11.1}{BSC 翻转率 $q$，输入概率 $(p,1-p)$。(a) 算互信息，证明 $C=1-\B(q)$。(b) 允许 5\% 错误概率时，求可达码率的最紧上界。}
\solve (a) 输出 0 来自输入 0 未翻转或输入 1 翻转，故 $p_Y(0)=p(1-q)+(1-p)q$。条件输出熵总为 $\B(q)$，所以 $I=\B(q+p(1-2q))-\B(q)$。首项至多 1，$p=1/2$ 时达到，故 $C=1-\B(q)$。(b) 按平均消息位错率解，共用推导给 $R\le[1-\B(q)]/[1-\B(.05)]$，分母约 $.713603$。若题中 5\% 指固定整块错误率，渐近上界由强逆定理给 $C$；若指信道翻转率则代入 $q=.05$，容量约 $.713603$。
\problem{11.2}{BEC 擦除率 $q$，输入概率 $(p,1-p)$。(a) 算互信息并证 $C=1-q$。(b,c) 容许错误概率 25\%、50\% 时分别求最紧码率上界。}
\solve (a) 收到 0 或 1 时输入确定；收到 $e$ 时，由于擦除与输入独立，后验不变。因此 $H(X|Y)=q\B(p)$，互信息 $(1-q)\B(p)$ 在 $p=1/2$ 取 $C=1-q$。(b) 平均消息位错率 $.25$ 给 $R\le(1-q)/[1-\B(.25)]\approx(1-q)/.188722$。(c) 允许 $.5$ 时，什么都不传，对均匀位盲猜的平均错误率就为 $.5$，故此条件不给有限码率上界。若百分比指固定整块错误率，两问的渐近上界都是 $C$；若指信道擦除率，容量分别是 $.75$、$.5$。
\problem{11.3}{输入 $a,b,c$，输出 $a,b,c,d$，转移矩阵四行是 $(3,1,0)/4$、$(1,2,1)/4$、$(0,1,2)/4$、$(0,0,1)/4$。输入概率 $(1/4,1/4,1/2)$。(a) 算互信息。(b) 对每个输出用最大 $P(y|x)$ 的输入译码，算平均错误率。(c,d) 码率分别为 $.01$ 和 $100$ bit/次、位误率 1\% 的声明可能吗？}
\solve (a) 矩阵乘输入：$p_Y(a)=3/16+1/16=1/4$；$p_Y(b)=1/16+2/16+1/8=5/16$；$p_Y(c)=1/16+1/4=5/16$；$p_Y(d)=1/8$。第一列非零概率为 $(3/4,1/4)$，后两列均有 $(1/2,1/4,1/4)$。故 $H(Y|X)=\tfrac14\B(1/4)+\tfrac34(1.5)$；$H(Y)=\tfrac14(2)-2(\tfrac5{16})\log(\tfrac5{16})+\tfrac18(3)$。相减得 $I\approx.59598$。
\parta{b} 按每行最大值，$a,b,c,d$ 分别译成 $a,b,c,c$。每种正确事件用先验乘该条件概率，正确率为 $3/16+1/8+1/4+1/8=11/16$，错率 $5/16$。
\parta{c,d} (a) 的输入已达到 $I>.01$，故 $C>.01$，足够长的码可使整块、从而平均位错误率都低于 1\%。另一方面 $C\le H(X)\le\log3$ 对最大化取值仍成立。共用推导给 $R\le\log3/[1-\B(.01)]\approx1.7243$；100 不可能。
\problem{11.4}{独立加性噪声信道 $Y=X+Z$，$X\in\{0,1\}$，噪声等概率取 $0,a$，$a$ 为任意整数。求容量。}
\solve 给定输入 0、1 时，输出集合分别为 $\{0,a\},\{1,1+a\}$。交集非空仅当 $a=1$（共享 1）或 $a=-1$（共享 0）。共享输出在两输入下都以 $1/2$ 出现，后验保持先验；其它输出确定输入。于是 $H(X|Y)=\tfrac12H(X)$，互信息最大为 $1/2$。其它 $a$ 时集合不交（$a=0$ 按噪声恒 0 理解），所以 $H(X|Y)=0$，均匀输入给 $C=1$。

\topic{Tut 12 · Hamming 与算术编码}
\problem{12.1}{对 $(7,4)$ Hamming 码，(a) 编码 \bb{0111}；(b) 用题目所指维恩图解码 \bb{1100000}；(c) 给该串的校验综合；(d) 最多保证纠正几位？同样 4 个数据位，要用多少位的重复码才能保证相同能力？}
\solve 原题 PDF 缺图，采用常见偶校验约定：数据位置 $3,5,6,7$，三组位置 $(1,3,5,7)$、$(2,3,6,7)$、$(4,5,6,7)$。
\parta{a} 数据位依次 $0,1,1,1$。第一组其余位异或 $0\oplus1\oplus1=0$，第二组 $0\oplus1\oplus1=0$，第三组 $1\oplus1\oplus1=1$，故校验位 $(p_1,p_2,p_4)=(0,0,1)$，组装得 \bb{0001111}。
\parta{b,c} 对 \bb{1100000} 按三组求实际奇偶，得到 $(s_1,s_2,s_4)=(1,1,0)$。综合按高位在前为 \bb{011}，即错误位置 $1+2=3$。翻第 3 位得 \bb{1110000}，取数据位得 \bb{1000}。
\parta{d} 校验矩阵的七列是不同非零的三位位置编号，单列或两列异或都不为零，所以非零码字至少三位；编号 $1,2,3$ 的列异或为零，距离可达 3。半径 1 的译码球互斥，半径 2 的可能重叠，保证纠一位。重复两次遇一错会平票，三次多数决可纠一错，发 $4\cdot3=12$ 位。
\problem{12.2}{首位 0、1 各半。之后下一个符号与前一位相同的概率 $3/5$，不同的概率 $1/5$，结束 EOS 概率 $1/5$。(a) 算术编码 \bb{011}；(b) 解回；(c) 自选一串编码并解回。}
\solve 按 $0,1,\code$ 排序，首轮只有 $0,1$。
\parta{a} 首位 $0$ 取 $[0,.5)$。前位为 0 时，$1$ 的累积段是 $[.6,.8)$，映到 $[.3,.4)$。前位成了 1，$1$ 的累积段是 $[.2,.8)$，映到 $[.32,.38)$。再取 EOS 的最后 20\%，区间 $[.32+.06(.8),.38)=[.368,.38)$。$[95/256,96/256)$ 在其中，所以发 95 的 8 位二进制 \bb{01011111}。
\parta{b} 代表值 $r=95/256=.37109375$。全区间 $r<.5$，读 0；在 $[0,.5)$ 中，$r$ 落在 1 段 $[.3,.4)$；该区间中落在 1 段 $[.32,.38)$；再落在 EOS 段 $[.368,.38)$，读出 \bb{011}。
\parta{c} 选 \bb{0}。首符号取 $[0,.5)$，EOS 切出 $[.4,.5)$。取 $[7/16,8/16)$，发 \bb{0111}。$r=7/16$ 先落 0 段，再落 $[.4,.5)$ 的 EOS 段，译为 \bb{0}。
\problem{12.3}{(a) 为何使用自适应分布？(b) Dirichlet 分布的直观含义？(c) 符号 $a,b,c$ 的虚拟计数为 $(3,1,1)$，EOS 固定 $1/4$。以自适应算术码编码 \bb{bb}。(d) 解回。}
\solve (a) 源的概率未知或会变化时，按观察更新，常见符号获得较宽的区间；编解码同步使用同一更新规则。(b) Dirichlet 是对概率向量的分布。参数可看作虚拟计数；每观察某符号，其计数加 1，预测概率是该计数占总计数的比例。
\parta{c} 选顺序 $a,b,c,\code$。前三项的计数比乘 $3/4$。开始 $a,b,c$ 的概率为 $.45,.15,.15$，故第一个 $b$ 取 $[.45,.60)$。更新 $(3,2,1)$ 后，$a,b$ 概率 $.375,.25$；第二个 $b$ 取 $[.45+.15(.375),.45+.15(.625))=[.50625,.54375)$。更新 $(3,3,1)$；EOS 仍是末 $1/4$，起点 $.50625+.0375(.75)=.534375$，终点 $.54375$。选 $[137/256,138/256)$，发 \bb{10001001}。
\parta{d} $r=137/256=.53515625$。首轮在 $[.45,.60)$，译 $b$ 并更新 $(3,2,1)$；第二轮在 $[.50625,.54375)$，再译 $b$，更新 $(3,3,1)$；第三轮在 $[.534375,.54375)$，读 EOS 停止，输出 \bb{bb}。
\end{multicols}
\end{document}
